
\documentclass{article} % For LaTeX2e
\usepackage[utf8]{inputenc}
\usepackage[T1]{fontenc}

\usepackage{times}
 \newcommand{\RNum}[1]{\uppercase\expandafter{\romannumeral #1\relax}}

% Optional math commands from https://github.com/goodfeli/dlbook_notation.
\input{math_commands.tex}

\usepackage{hyperref}
\usepackage{url}
\usepackage{bbm}  %For indicator
\usepackage{subcaption}
\usepackage{xcolor}
% profo environment
% \renewcommand{\qedsymbol}{$\blacksquare$}
\usepackage{amsthm}


% Define the \TK command
\newcommand{\TK}[1]{\textcolor{red}{TK: #1}}

\usepackage{comment}
\newcount\Comments  % 0 suppresses notes to selves in text
\Comments = 1
\newcommand{\kibitz}[2]{\ifnum\Comments=1{\color{#1}{#2}}\fi}
\definecolor{english}{rgb}{0.0, 0.5, 0.0}
\newcommand{\dcpadd}[1]{\kibitz{english}{#1}}
\newcommand{\dcp}[1]{\kibitz{teal}{[DCP: #1]}}
\newcommand{\jeffrey}[1]{\kibitz{red}{[jeffrey: #1]}}
\newcommand{\kerem}[1]{\kibitz{cyan}{[kerem: #1]}}
\newcommand{\jeff}[1]{\kibitz{brown}{[jeff: #1]}}
\newcommand{\avs}[1]{\kibitz{purple}{[avs: #1]}}
\newcommand{\kz}[1]{\kibitz{blue}{[kz: #1]}}


% Authors must not appear in the submitted version. They should be hidden
% as long as the \iclrfinalcopy macro remains commented out below.
% Non-anonymous submissions will be rejected without review.

% \documentclass[12pt]{article}
\usepackage{amsmath} 
\usepackage{graphicx}
\usepackage{cite} 
\usepackage{hyperref}
\usepackage{natbib}
\usepackage{float}
\usepackage{geometry}
\usepackage{pgfplots}
\usepackage[framemethod=TikZ]{mdframed}
\usepackage{enumitem}
\usepackage{lipsum}
\usepackage{subcaption}
\usepackage[most]{tcolorbox}
\usepackage{setspace}
\usepackage{textcomp}
\usepackage{tabularx}
\usepackage{array}
\usepackage{tikz}
\usepackage{multirow}
\usetikzlibrary{arrows,positioning, automata,arrows.meta, shapes.geometric, decorations.pathreplacing}
\usepackage{booktabs} % For intervention table offset
% \input{parameters/params.tex}
\captionsetup{justification=raggedright,singlelinecheck=false}

% 3 -- WORDING & STYLING
\newcommand{\COMMENT}[1]{\textbf{\textcolor{red}{#1}}}
\title{\vspace*{-1cm} Learning from Synthetic Labs: \\ Language Models as Auction Participants \thanks{
  \footnotesize
This research was made possible by a generous grant from Dropbox Inc. Thanks to Parker Whitfill and Kobi Gal for their helpful feedback. Author's contact information, code, and data will be available at https://github.com/KeHang-Zhu/llm-auction. Both Anand V. Shah and Kehang Zhu contributed equally to this work. John J. Horton is a co-founder of Expected Parrot Inc., which develops generative AI models for market research applications. The views expressed herein are those of the authors and do not necessarily reflect the views of the National Bureau of Economic Research.
}
}

  
\author{
  Anand Shah$^{*}$\\
  MIT \\
  \and
  Kehang Zhu$^{*}$\\
  Harvard  \\
  \and
  Yanchen Jiang \\
  Harvard \\
  \and
  Jeffrey G. Wang \\
  Harvard \\
  \and
  Arif K. Dayi \\
  Harvard \\
  \and
  John J. Horton \\
  MIT \& NBER \\
  \and
  David C. Parkes \\
  Harvard \\
}
\date{\today} 


% \author{Anand, Kehang \& Yancheng Jiang\thanks{ Use footnote for providing further information
% about author (webpage, alternative address) -- \emph{not} for acknowledging
% funding agencies.  Funding acknowledgements go at the end of the paper.} \\
% Department of Computer Science\\
% Cranberry-Lemon University\\
% Pittsburgh, PA 15213, USA \\
% \texttt{\{hippo,brain,jen\}@cs.cranberry-lemon.edu} \\
% \And
% Ji Q. Ren \& Yevgeny LeNet \\
% Department of Computational Neuroscience \\
% University of the Witwatersrand \\
% Joburg, South Africa \\
% \texttt{\{robot,net\}@wits.ac.za} \\
% \AND
% Coauthor \\
% Affiliation \\
% Address \\
% \texttt{email}
% }

% The \author macro works with any number of authors. There are two commands
% used to separate the names and addresses of multiple authors: \And and \AND.
%
% Using \And between authors leaves it to \LaTeX{} to determine where to break
% the lines. Using \AND forces a linebreak at that point. So, if \LaTeX{}
% puts 3 of 4 authors names on the first line, and the last on the second
% line, try using \AND instead of \And before the third author name.

\newcommand{\fix}{\marginpar{FIX}}
\newcommand{\new}{\marginpar{NEW}}

%\iclrfinalcopy % Uncomment for camera-ready version, but NOT for submission.
\begin{document}

\maketitle
\begin{abstract}
This paper investigates the behavior of simulated AI agents (large language models, or LLMs) in auctions, introducing a novel synthetic data-generating process to help facilitate the study and design of auctions. We find that LLMs -- when endowed with chain of thought reasoning capacity -- agree with the experimental literature in auctions across a variety of classic auction formats. In particular, we find that LLM bidders produce results consistent with risk-averse human bidders; that they perform closer to theoretical predictions in obviously strategy-proof auctions; and, that they succumb to the winner's curse in common value settings. On prompting, we find that LLMs are not very sensitive to naive changes in prompts (e.g., language, currency) but can improve dramatically towards theoretical predictions with the right mental model (i.e., the language of Nash deviations). We run 1,000$+$ auctions for less than $\$$400 with GPT-4 models (three orders of magnitude cheaper than modern auction experiments) and develop a framework flexible enough to run auction experiments with any LLM model and a wide range of auction design specifications, facilitating further experimental study by decreasing costs and serving as a proof-of-concept for the use of LLM proxies. 
\end{abstract}

\begin{titlepage}

\maketitle\makeatletter \gdef\@ACM@checkaffil{} \makeatother
% Optionally include a table of contents
% \vspace{-0.1cm}
\setcounter{tocdepth}{2} % adjust to 1 if desired
% \tableofcontents

\end{titlepage}
\newpage

% \tableofcontents

\newpage
\section{Introduction} \label{sec:introduction}
% The field of mechanism design, auction theory in particular, has benefited enormously from a rich interplay between empirics and theory. One recent example is the development of {\em obviously strategy-proof mechanisms} (hereafter, OSP) \citep{li2017obviously}. This technical refinement was inspired by an empirical puzzle: despite it being well-known that open-ascending clock and second-price sealed-bid auctions were strategically equivalent in private value settings,
% experiments since the 1980s suggested that people are much better at playing the clock format than the sealed-bid format ~\citep{kagel1987information}. Motivated by empirical evidence, OSP gives an explanation for why the clock format might be easier to play than a sealed-bid format and has since inspired a groundswell of work in auction design under behavioral constraints. The story echoes a well-understood but worth emphasizing point: empirical work is vital to the development of new theory. 

% Yet, acquiring the right empirical can be extremely expensive. Even \citet{li2017obviously}'s relatively contained laboratory experiment cost over \$15{,}000, employed over 400 participants, and likely cost hundreds of researcher-hours to set up.
% These types of controlled experiments are also the best case scenario. The

Rapidly advancing capabilities of large language models (LLMs) raise an exciting possibility: might LLM agents serve as cost-effective proxies for human bidders, thereby making experiments orders of magnitude cheaper to run? In this paper, we contribute to an emerging literature on LLM proxies by testing whether LLM agents can replicate key empirical regularities of the auctions literature. We map LLM behavior in standard auction formats against established human and theoretical benchmarks from economics, providing a proof-of-concept for low-cost, large-scale auction experimentation using LLM bidders.

We begin by benchmarking play with LLM bidders in classic environments: sealed-bid auctions in independent, private-value settings. When there are departures from the theory, we are interested in whether the departures agree with empirical results \citep{kagel1993independent, kagel2020handbook}. We find that agents fail to replicate revenue equivalence results (as existing empirical evidence suggests). We find that bids in the SPSB auction are higher than bids under the FPSB auction, as would be expected, but that there is a smaller separation between the two than predicted by theory. This is primarily due to bids under the FPSB auction being higher than the risk-neutral Bayes-Nash equilibrium suggests. One possible explanation is that LLMs play according to some level of risk aversion: this is consistent with \citet{cox1988theory}'s survey of over 1,500 IPV auction experiments, which finds that revenue under the FPSB auction is higher than under the SPSB auction and suggests that this is due to risk aversion. In addition, we see LLMs in the SPSB auction tend to submit bids lower than their value to an extent that is not echoed in the experimental economics literature. We also supplement this analysis with simulations of the third-price sealed bid (TPSB) and all-pay auctions.

Next, we fix the second-price auction and consider classic ways to make the auction `easier' or `harder' to play, to learn whether LLMs exhibit behavior similar to humans in the face of cognitive constraints. 

Theoretical and experimental studies~\citep{li2017obviously} argue that ascending clock formats are less cognitively demanding than sealed-bid ones, causing humans to make less mistakes relative to the dominant strategy of bidding one's value. We find that LLMs also find clock formats easier to play -- they are more likely to exit the auction (the analogue to bidding in the sealed-bid format) at their value. Following \citet{li2017obviously}, we study this in an affiliated private value setting, replicating their results with LLMs. We also compare the typical clock auction with a blind variant (a clock auction where bidders do not know when others drop out) and the corresponding second-price sealed-bid auction, and find that switching to a clock format accounts for most of the improvement in LLM play. For robustness, we also run these experiments in an IPV setting and obtain the same results.


Next, we add complexity to the auction environment by considering the second-price auction under a common value setting. Under common values, all agents agree on their value for the prize $c$, but each only observes their valuation for the prize under some idiosyncratic, private measurement error $c + p_i$. This auction environment is more cognitively demanding than the independent, private values case as agents do not know whether their observed high valuations are due to high draws of $c$ or the measurement error $p_i$, and this tension produces the winner's curse: that the winning agents wish they had never won, having been adversely selected to be the bidder who least appreciated their high signal was due to a high $p_i$. The winner's curse is purely a non-rational phenomenon -- only the cognitively limited agent suffers from it. The existing economics literature finds humans regularly suffer the winner's curse in common value settings, and we find results vindicating this evidence with LLM agents: LLMs succumb to the winner's curse as well. In particular, \citet{kagel1986winner} famously find experimental evidence for the winner's curse and argue that the winner's curse barely bites with 3-4 bidders but seriously bites for larger auctions with 6+ bidders. We find qualitatively similar results hold with LLM agents as well. 

% \avs{need to figure out learning (there must be \textit{some } conditions under which learning holds. what are they?) if we want this to be a major result.}

These classic results are some of the reasons economists have found auction theory to be beautiful over the last few decades, but auctions can also be useful \citep{milgrom2020clock}. Auctions are used to clear billions of dollars of value annually, and one may ask how LLMs fare in less controlled settings. To answer this, we design a more realistic auction environment inspired by eBay's online marketplace. We examine two design parameters: a hidden reserve price and a modified closing rule (soft-close, which extends the auction when there is bidding activity). First, we demonstrate that LLM agents replicate emergent, real-world bidding behaviors seen on eBay with no prompting: in auctions with a hard closing time, LLMs frequently delay bids until the last second to `snipe' their competitors as humans would. Second, experiments testing the design of an extended closing rule greatly reduce bid sniping and improve price discovery, echoing dynamics of online auctions between Amazon and eBay in the early 2000s \citep{roth2002last}.

The ability to interface more directly with our experimental subjects also enables us to study the explanation of mechanisms more carefully. Auctions are most prized as mechanisms to allocate goods efficiently, but if agents don't understand the rules they can't be expected to play well. How can we describe and explain mechanisms to empower optimal play cheaply? We conclude the body of our paper by testing interventions designed to influence or aid LLM bidders' understanding of the underlying economics of the setting and its willingness to follow rules. In particular, we test six interventions: an extensive-form menu description meant to simplify the equilibrium logic inspired by \citet{gonczarowski2023strategyproofness}; a description describing the auction via a clock (without actually running the clock auction) inspired by \citet{breitmoser2022obviousness}; a description walking an LLM through the logic of Nash deviations; a description which also reveals the correct statement that the SPSB's dominant strategy is to bid one's value; and a description that incorrectly reveals a false statement that the SPSB's dominant strategy is to bid half of one's value. We find that that all the interventions that were meant to improve the correctness of an LLM agent's bidding strategy helped, with the logic of Nash deviations helping the most of the first three. Unsurprisingly, the false statement intervention caused LLMs to play poorly, but interestingly, the LLM continued to experiment with new strategies upon determining the given advice was bad.

The remainder of the paper proceeds as follows. In Section 1.1, we report related work. In Section 2, we describe our simulation procedure and prompting methods in more depth. In Section \ref{section:classic}, we benchmark LLM play against sealed-bid auctions in independent, private value settings. Then, for the rest of the paper, we fix the game as a second-price auction and test various design considerations. In Section \ref{session:OSP}, we vary the auction format and find that LLMs play clock formats the best. In Section \ref{session:winner}, we add cognitive complexity to the auction environment by switching to a common value environment and show that LLMs succumb to the winner's curse. Both of these sections replicate key empirical regularities found in laboratory and real-life play in auctions with humans. In Section~\ref{session:ebay}, we implement auctions inspired by eBay's marketplace and provide new evidence on famous debates on design parameters. In Section~\ref{session:Intervention}, we introduce interventions designed to influence or aid LLM bidders' understanding of the underlying economics of the setting. We find that prompting matters, and that, surprisingly, LLMs improve play dramatically when instructed to follow the logic of Nash deviations.

In the Appendix we report complete simulation procedures, prompts and numerous robustness checks. In particular, we ran experiments with prompts from other contemporaneous auctions experiments (e.g., the Appendix script from \citet{li2017obviously}), with prompts in different languages and with different currencies, and with varying numbers of bidders. All experiments are run with chain-of-thought agents, but we also run an ablation study with out-of-the-box agents. We also find interesting evidence that `goal' prompting (e.g., reminding the LLM that the goal is to maximize profit) leads to bidding that more closely follows rational economic theory, leading to higher allocative efficiency. See also~\citet{manning2024automated} for a related discussion. 

To obtain the data for these empirical results, we have developed a code repository to systematically run experiments 
with a certain number of bidders and a particular prompt. 
In particular, our repository is flexible enough that it can be used to generate synthetic data for almost any describable  auction format and including auctions with single or multiple goods.\footnote{We will make the code-base public  and hope this will facilitate additional research into the role of LLMs as modeling human bidder behavior.}
%
For the experiments described in this paper, we ran more than $1,000$ auctions with more than $5,000$ GPT-4 agent participants for costs totaling less than $\$400$. 
In contrast, the largest survey of auction experiments to date comes from \citet{cox1988theory}, reporting the results of
running $1,500+$ auctions and with considerably higher total cost.


% Head 1
\subsection{Related Work}




\textbf{Auctions:}
There's a vast quantity of theoretical and experimental literature on auctions. 
While \citet{krishna2009auction}'s textbook and \citet{kagel2020handbook}'s handbook provided invaluable general resources, we will stay focused only on the citations relevant to the results in this paper.

Starting with the IPV case, the benchmark of revenue equivalence in the risk-neutral case is exposited in the seminal 
\citet{myerson1981optimal}. 
The experimental evidence for departures due to risk-aversion in the FPSB auction is thoroughly documented in \citet{coppinger1980incentives} and \citet{cox1988theory}'s survey. 
The experimental evidence for the common error of bidding above one's value in the SPSB auction is documented in \citet{kagel1993independent}.

For common value settings, the winner's curse has been documented empirically and experimentally. 
Experimental evidence for the winner's curse was first given by \citet{bazerman1983won}. 
For the present paper, we primarily follow the empirical approach in \citet{kagel1986winner} (FPSB auctions with varying numbers of bidders) and \citet{kagel1987information} (SPSB auctions with informational interventions). 
Empirically, the winner's curse was first documented in 1971, in auctions for oil drilling rights by \citet{capen1971competitive}. 

Theoretically, while the winner's curse is fundamentally a problem of non-rational play,  \citet{charness2009origin} provide a model to cast the adverse-selection problem as a failure from cognitive constraints. 


Recent work on obvious strategy-proofness began with \citet{li2017obviously}, who demonstrates empirically that human subjects tend to be less truthful in second price sealed bid auctions than ascending clock auctions in the APV setting, even though the two auctions are strategically equivalent. 
Li also provides a theoretical framework for the results. 
To better understand our simulations, we also consider the experimental evidence presented by \citet{breitmoser2022obviousness}, who investigate intermediate auction formats that decompose the behavioral effects in \citet{li2017obviously}. 

\textbf{LLMs as simulated agents:} Recent LLMs, having been trained on an enormous corpus of human-generated data, are able to generate human-like text and reasoning \citep{achiam2023gpt, bubeck2023sparks}. 
Yet, they are far from perfect -- in particular, displaying limited planning abilities and reflecting various cognitive biases endemic to human agents \citep{wan2023kelly}. 

There is a growing literature on using these human-like AI models as simulated agents in economics and social science studies \citep{aher2023using, park2023generative, brand2023using}. 
In this literature, \citet{horton2023large} replicates four classical behavioral economics experiments by endowing a single LLM agent with different personas and querying it about its decisions and preferences.
\citet{manning2024automated} enable multiple GPT-4 agents to interact and simulate various social science scenarios, including bargaining, job interviews, and bail-setting.
Finally, \citet{raman2024steerassessingeconomicrationality} benchmark the ability of LLM agents to conduct rational play over a broad range of tasks.

Compared to humans, the inference cost of LLM queries is very low and continues to decrease \citep{achiam2023gpt,patel2023splitwise,bae2023complexitynet}.

Beyond acting as decision-makers, LLMs can also prepare the very environments in which simulated agents operate.
Recent work shows that LLMs can parse natural-language narratives (e.g., novels, news articles, or other free form texts) and translate them into formal game-theoretic settings, thereby generating the strategic context required for agent-based experiments \citep{daskalakis2024charting, mensfelt2024autoformalization, deng2025natural}.
In parallel, LLMs are being used as preference-elicitation tools that help existing economic mechanisms gather richer, textual inputs from humans \citep{soumalias2025llm, huang2025accelerated}.


\textbf{LLMs in auctions:}
There are some works that use LLMs as simulated agents in auction experiments.
\citet{fish2024algorithmic} study the collusion behaviors in first-price sealed-bid auction of two LLM agents
 under the context of LLMs as a price setter for companies. 
\citet{chen2023put} study how to make an LLM better at playing auctions than humans.
Finally, \citet{manning2024automated} run a more limited study an a variant of an open-ascending clock auction with three bidders, focusing on deviations from rational economic theory in considering bidders' values and the final clearing price.
% \dcpadd{ran a more limited study an  a variant of an open-ascending clock auction with three bidders, 
% focusing on deviations from rational economic theory in considering bidders' values and 
% the final clearing price.}

\textbf{Automated Mechanism Design:} Pioneering work by \citet{conitzer2002complexity, conitzer2004self} introduced \emph{automated mechanism design} (AMD), which uses computational search to design optimal, incentive-compatible auctions \citep{curry102automated}.
A newer line -- often labeled “differentiable economics’’ \citep{dutting2024optimal} -- pushes this idea further with deep learning, producing optimal auctions \citep{dutting2024optimal}, combinatorial auctions \citep{wang2025bundleflow}, budget-constrained bidders \citep{feng2018deep}, payment minimization \citep{tacchetti2019neural}, two-sided matching \citep{ravindranath2021deep}, and data-market pricing \citep{ravindranathdata}.
Structural advances such as \citet{wang2024gemnet} obtain \emph{exact} strategy-proofness rather than approximate incentive compatibility. Whereas automated-mechanism-design (AMD) research pushes the analytical frontier of optimal auction design, our study advances the behavioral understanding of how LLM-based agents act within classical formats. 


\section{Methods}

\subsection{LLM as Auction Participants}

We use a single LLM API call to represent one bidder in our auction simulations. Unlike traditional agent-based modeling, we do not anchor bidder behaviors or personas mechanically beyond a name to index them (e.g., Alice, Betty or Charles). Instead, we allow LLMs to conduct their own strategic reasoning and in-context learning to generate bids. Hence, we focus on how these LLM agents strategically interact, without concern for demographic representativeness, as our primary interest is rational and strategic behavior in the mechanism.

To ensure that LLM agents make profit-maximizing decisions, we append a prompt prefix that explicitly instructs them to seek long-run profit maximization. This instruction is universal in all the experiments ran in this paper, and is inspired by \citep{fish2024algorithmic, chen2023put}. 
Complete details on prompts and execution can be found in Appendix ~\ref{app:process}. Specific implementation details for non-sealed-bid formats (the clock and eBay auction formats) are provided in their respective sections.

Before eliciting bids in each period, we provide a description of the auction rules along with randomly drawn values and information from previous rounds. Standard prompting techniques like chain-of-thought \citep{wei2022chain} are also used to mimic the strategic thinking period in auction practices \citep{li2017obviously}.

% \begin{quote}
% \begin{lstlisting}[basicstyle=\ttfamily\small]
% Your TOP PRIORITY is to place bids which maximize your 
% profit in the long run. To do this, you should explore 
% many different bidding strategies, including possibly 
% risky or aggressive options for data-gathering purposes. 
% Learn from the history of previous rounds in order to 
% maximize your total profit. Don't forget the values are 
% redrawn independently each round.
% \end{lstlisting}
% \end{quote}

% We observed that LLM agents become more effective in auctions once told explicitly to profit-maximize. This motivation encourages the agents to adapt and try various strategies for maximizing total earnings.


\subsection{Simulation Procedure}\label{sec:process}

For each auction experiment, we simulate $n$ (often 3) LLM agents bidding against one another. Each setting is repeated multiple times, with randomized private values drawn for each round. We design our multi-round protocol to mirror standard laboratory experiments in the auction literature. The item to be bid in standard auctions is specified as a ``prize''. In eBay auctions, we arbitrarily set the prize item as a ``256GB IPhone 16 pro'' and the item condition is set to be ``used'' for realism.

Below is our simulation procedure for each auction; full prompt texts are in Appendix~\ref{app:prompt}, and an example round is in Appendix~\ref{app:process}:
\begin{enumerate}
    \item {\textit{Planning.} LLM agents receive a description of the auction format and are asked to propose a bidding plan.}
    \item {\textit{Value Realization.} Each LLM agent is assigned a random private value.}
    \item {\textit{Bidding.} Given their value, each agent places a bid according to their stated plan.}
    \item {\textit{Outcome.} The system computes allocation and payments. Each LLM agent learns its own profit or loss outcome, as well as the full set of bids submitted (but not other agents' private values).}
    \item {\textit{Reflection.} Before the next round, each LLM agent reviews the results and possibly revises its approach. These updates are stored in the \texttt{HISTORY} variable.}
    \item {\textit{Repetition.} One complete “plan-bid-reflect” loop is one round. We repeat 15 rounds in each experiment, with the entire \texttt{HISTORY} carried forward.}
\end{enumerate}

API calls are made using the Python package Expected Parrot Domain-Specific Language (EDSL), a wrapper that structures LLM queries and responses in the desired format \citep{Horton2024EDSL}. Each agent in our sample is supported by a separate API call via EDSL, thereby nullifying the risk of collusion. For every auction, we use models from OpenAI at a temperature of $0.5$.\footnote{{We set a non-zero temperature to encourage diversity in strategies and reflections. No parameter tuning was performed to achieve these results.}} 
The full code base is made open-source for research and replication at \href{https://github.com/KeHang-Zhu/llm-auction}{GitHub}.
% \avs{We make this live when the paper is ready for NBER.}

\subsection{Discrete Bids and Continuous-Theory Benchmarks}

A fundamental tension in experimental auctions work lies between the use of continuous theory and the necessarily discrete nature of implementation. Auction theory traditionally models bids and valuations as continuous variables, enabling the use of calculus to derive clean analytical results. Yet  real-world auctions -- whether in experimental laboratories or practical markets -- are invariably discrete, with quantized valuations (e.g., integer values from 0 to 99) and finite bid increments (e.g., \$1). This discrete-continuous gap has received surprisingly limited direct attention in the literature, raising the critical question of how to justify the use of continuous-theory benchmarks for our discrete simulations. 

The foundational literature in both theoretical and experimental economics provides a basis for this practice. A deep and influential body of work operates on the principle that, with reasonably fine increments, discrete auctions serve as a good approximation for the continuous counterparts. In many parts of the literature, this intuition is almost implicit (e.g., as in footnote 3 of \cite{noussair2006behavior}), but it has also been formalized in other works, as in \cite{cai2010note}'s convergence results that outcomes in discrete first-price auctions converge to the continuous-case equilibrium as bid increments shrink. This theoretical backing underpins the standard methodology in experimental economics, where pioneers like \cite{kagel1993independent} have long benchmarked behavior in discrete laboratory settings against the predictions of continuous theory.

However, a more modern line of inquiry reveals that convergence arguments do not generalize as smoothly as once hoped, and that model equilibria can be sensitive to environment parameters (e.g., tie-breaking rules, number of bidders, and auction type). This is all the more striking given foundational results like \cite{athey2001single} provide quite general conditions -- most notably, single crossing -- that guarantee the existence of a pure-strategy equilibrium. Yet, even when such equilibria exist, their specific \textit{form} may be fragile. Building on early work from \cite{maskin1985revenue}, recent ``discontinuity'' results from \cite{rasooly2020importance} show that familiar, symmetric pure-strategy equilibria often vanish when bidder valuations and bids are restricted to a discrete grid, yielding only more complex asymmetric or mixed-strategy equilibria. This fragility is particularly acute for the all-pay auction, where Rasooly \& Gavidia-Calderon show that a symmetric, pure-strategy equilibrium fails to exist for a grid of values and bids (no matter how fine), regardless of the tie-breaking rule. This unconditional non-existence is stronger than for the first-price auction, where pure, symmetric equilibria can still exist for a low number of bidders and an even number of possible values. And, this fragility is not just destructive; an example from \cite{boudreau2011all} shows that in some contests, discreteness can be constructive, creating pure-strategy equilibria where none existed in the continuous model.

This leaves the relationship between discrete and continuous auction models as a live and important frontier of research. If we are to continue leveraging the powerful insights from classic auction theory, we must better understand the bridge between discrete and continuous models. The methods of this paper, insofar as they dramatically reduce the cost of obtaining synthetic evidence on auctions, enables progress towards that goal. In particular, future work can feasibly investigate how outcomes change as we vary parameters that theory tells us are operative -- the number of bidders, the granularity of the bid grid, and the specific auction format -- in a way that is currently prohibitively expensive with human subjects.\footnote{In particular, our codebase on GitHub implements a parameter of minimal increment to allow researchers to explore the effect of bidding increments on auctions.} This computational advantage could, in future work, help identify the boundary conditions where continuous models remain valid and reveal strategic phenomena unique to discrete settings.
% \dcp{do you show these results, exploring discretization effects? if so, point to them here. if not, drop the remark.}

% positions LLM-based auction experiments to contribute not just to our understanding of mechanism design, but also to the  theoretical question as to  when and how continuous models accurately represent discrete strategic interactions. By generating large datasets across varied discrete specifications, synthetic experiments can help identify the boundary conditions where continuous approximations remain valid and reveal the strategic phenomena that emerge uniquely in discrete settings. In this way, our methodology offers a path forward for reconciling the mathematical elegance of continuous auction theory with the practical constraints of discrete bids. \dcp{is this over-promising for the current paper? best shortened, and situated as "In the future, this computational advantage of .. could ... "?}


\section{Benchmarking Against Existing Auctions}\label{results}


\subsection{Sealed Bid Auctions with Independent Private Values}\label{section:classic}
We begin by examining the First-Price Sealed-Bid (FPSB), Second-Price Sealed-Bid (SPSB), Third-Price Sealed-Bid (TPSB) and All-Pay Sealed-Bid auctions in an independent private values (IPV) setting. 


\subsubsection{Setting}

There are three bidders in each auction, and each bidder $i$'s value is drawn from an independent, uniform distribution $v_i \sim U[0, 99]$. Bidders, upon observing their value, submit a sealed bid $\beta(v)$, with $\beta(\cdot)$ a vector mapping each component value to its corresponding bid, corresponding to a strategy profile for reporting based on the valuation profile. 

In the FPSB auction, the highest bidder pays her bid and receives the prize (and all other bidders pay $0$ and receive no prize). Formally, the payment for an agent is given by: $$t_i(\beta(v)) = \mathbbm{1}_{i \; \text{won the auction}}\cdot \beta_i(v).$$
In the SPSB auction, the highest bidder pays the second-highest bid and receives the prize (and all other bidders pay $0$ and receive no prize). 
The payment for an agent is given by: $$t_i(\beta(v)) = \mathbbm{1}_{i \; \text{won the auction}} \cdot \beta^{(2)}(v),$$ where $\beta^{(2)}$ represents the second-order statistic or the second-largest bid. 
Analogously, in the TPSB auction, the highest bidder pays the third-highest bid and receives the prize (and all other bidders pay $0$ and receive no prize). The payment for an agent is given by: $$t_i(\beta(v)) = \mathbbm{1}_{i \; \text{won the auction}} \cdot \beta^{(3)}(v),$$ where $\beta^{(3)}$ represents the third-order statistic or the third-largest bid.
Finally, in the all-pay auction, all bidders pay their bid but only the highest bidder wins the prize. The payment for an agent in the all-pay auction is given by: $$t_i(\beta(v)) = \beta_i(v).$$
Bids are submitted in $\$1$ increments and ties are resolved randomly. 

\subsubsection{Theoretical benchmarks}\label{sec:ipv-theory}
The risk-neutral equilibrium strategy in all four of these auction formats is well understood. We summarize the functional forms of these strategies - specific to our setting - below in Table \ref{tab:risk-neutral-eqm}. 
\begin{table}[h]
    \centering
    \renewcommand{\arraystretch}{1.4} % Increases row spacing
    \begin{tabular}{l c}
        \hline
        \textbf{Auction Type} & \textbf{Risk-Neutral Equilibrium Strategy} \\ 
        \hline
        First-Price Sealed-Bid$^*$ & $\frac{2}{3}v_i$ \\ 
        Second-Price Sealed-Bid$^{**}$ & $v_i$ \\ 
        Third-Price Sealed-Bid$^*$ & $2v_i$ \\ 
        All-Pay Auction$^*$ & $\frac{2}{3 \cdot 99^2} v_i^3$ \\ 
        \vspace{-5mm} \\ % for a bit more padding on the last row
        \hline
    \end{tabular}
    \caption{Risk-neutral equilibria of IPV auctions with 3 bidders and values drawn uniformly, $v_i \sim U[0, 99].$ A single asterisk (*) denotes a Bayes-Nash equilibrium, while a double asterisk (**) denotes a dominant strategy equilibrium.}
    % Derivations in Appendix \ref{app:risk-neutral-equilibria-proofs}.
    \label{tab:risk-neutral-eqm}
\end{table}

% Risk-neutrality is a particularly strong theoretical assumption. In practice, human subjects deviate from the risk-neutral equilibria of Table \ref{tab:risk-neutral-eqm}, and these deviations have motivated a broader theoretical literature seeking to rationalize bidding behavior under alternative assumptions about bidder preferences and decision-making (ambiguity aversion, loss aversion, probability misperception, etc). The most enduring and broadly applicable explanation remains risk aversion -- we will focus on this explanation in the following Section as well. 

% Risk-averse bidders are willing to lose value in expectation for more certainty. In this paper, we focus on two canonical formalizations: Constant Absolute Risk Aversion (CARA) and Constant Relative Risk Aversion (CRRA) preferences.
% % CARA and CRRA Utility Functions
% \begin{align*}
%     \text{CARA:}& \quad U(x) = -e^{-\lambda x}, \quad \lambda > 0\\
%     \text{CRRA:}& \quad U(x) = \frac{x^{1-\rho}}{1-\rho}, \quad \rho > 0, \quad \rho \neq 1
% \end{align*}

% The CARA specification assumes a bidder's absolute risk tolerance remains fixed across wealth levels, whereas CRRA assumes risk tolerance scales with wealth. CARA is more analytically tractable (and has a closed form for each auction) but is not scale invariant (switching units from dollars to cents changes CARA parameter estimates). CRRA is less analytically tractable (we report numerical solutions for empirical results) and scale invariant, making it easier to compare our results with previous empirical results. 

% Table \ref{tab:risk-averse-eqm} presents the equilibrium bidding strategies under CARA and CRRA preferences, with functional forms that explicitly incorporate the risk-aversion parameter. There is no new theoretical prediction for the SPSB auction under risk-aversion, as dominant strategies are invariant to risk attitudes. 

% \begin{table}[h]
%     \centering
%     \renewcommand{\arraystretch}{1.6}
%     \begin{tabular}{l c c}
%         \hline 
%         \textbf{Auction Type} & \textbf{CARA Equilibrium} & \textbf{CRRA Equilibrium}\\ 
%         \hline\vspace{-5mm}\\ %\vspace for padding top row to make it look nicer!
%         First-Price Sealed-Bid & $v - \frac{1}{\lambda} \ln\left(\frac{n-1}{n} + \lambda v\right)$ & $v - \frac{1}{\rho} \left( \frac{(n-1)v}{n} \right)^{1-\rho}$ \\
%         Third-Price Sealed-Bid & $v + \frac{1}{\lambda} \ln\left(1 + \frac{\lambda v}{n-2}\right)$ & \avs{solve proof. no result exists yet.} \\
%         All-Pay Auction & $v - \frac{1}{\lambda} \ln\left( 1 + \lambda v \right)$ & \text{Numerical solution required} \\ 
%         \hline
%     \end{tabular}
%     \caption{Bayes-Nash equilibria under CARA ($\lambda$) and CRRA ($\rho$) preferences with $n$ bidders and values drawn uniformly, $v \sim U[0, 1].$ Derivations in Appendix \ref{app:risk-averse-equilibria-proofs}. \avs{These are not all correct. come back and fix.}}
%     \label{tab:risk-averse-eqm}
% \end{table}

% Finally, irrespective of risk attitudes, theory universally predicts monotone bidding functions in single-unit IPV settings. Monotone bidding is a benchmark of auction theory more generally, reflecting the assumption in utility theory that ``more is better''.

Finally, note that in contrast with the dominant strategy equilibrium in SPSB auctions, the BNE reported here are sensitive to the ex-ante symmetry in all of our auction experiments and to the
assumption of risk-neutrality. We exert no control over LLM agents’ risk preferences in these
simulations except in one intervention in Section \ref{session:Intervention}.

\subsubsection{Empirical benchmarks}

Results from laboratory experiments offer additional empirical benchmarks against which we can compare LLM performance. In general, empirical evidence suggests
monotone bidding strategies\cite{kagel1993independent}, \cite{kagel2020handbook}. This means that higher
valuations consistently lead agents to report higher bids. However, the specific bidding patterns
observed in experiments often deviate from theoretical equilibria in interesting ways.
\paragraph{Empirical findings for FPSB:} In the case of FPSB auctions, experimental evidence consistently
shows bids above the risk-neutral BNE prediction \cite{cox1988theory}. They identify this with bidder
risk-aversion. In an FPSB auction, bidding higher increases the probability of winning but decreases
the potential profit if one does win. Risk-averse bidders may be willing to accept lower potential
profits in exchange for a higher chance of winning, leading to bids above the risk-neutral BNE.
\paragraph{Empirical findings for SPSB:} Similarly, experimental data for SPSB auctions also reveals
evidence of overbidding.  \cite{kagel1987information} and \cite{kagel1993independent} provide insights into this
behavior, in particular highlighting that a significant proportion of participants (typically around
60-70\%) in SPSB auctions submitted bids above their true values. However, in this case, the SPSB
has an equilibrium in dominant strategies, so the explanation is not risk-aversion. These findings
highlight the complexity of human behavior in auction settings. Several explanations have been
proposed, including misconceptions about the second-price rule [\cite{kagel1987information}] and strategic
uncertainty or beliefs about other bidders’ irrationality [\cite{crawford2007level}
].
\paragraph{Empirical findings for TPSB:} There is less empirical evidence for the TPSB. \cite{kagel1987information} test TPSBs in the laboratory setting and found that the empirically observed bidding strategy
was higher than both that of the SPSB and FPSB auctions, as theory predicts. In particular, 90% of
bids exceeded their value in this auction, and risk-aversion induces lower bidding (as opposed to
higher bidding in the case of the FPSB auction).
\paragraph{Empirical findings for all-pay auctions:} There is a plethora of evidence for all-pay auctions
or similar formats [\cite{dechenaux2015survey}]. As in previous formats, one major feature of bidding
behavior is overbidding relative to the BNE prediction: this holds in asymmetric and symmetric
auction environments and with various levels of information. Second, risk-aversion appears to push
bidding downwards (but still above the predicted level). Some experiments also report substantial
heterogeneity in bidding strategies, with bidding strategies differing sharply based on bidder types
(e.g., \cite{gneezy2006all}).

% For each of the IPV auction formats, we perform 5 experiments with 15 rounds and 3 LLM bidders competing over a prize.\footnote{While these specifications are designed to meet the `classic' auction formats of the literature, we also report the results of several variants in Appendix \ref{app:robustness}. In particular, we run one-shot auction simulations (as opposed to multi-round) with $n = 4, 5,$ and $10$ bidders, auctions with prompts using various currencies, and auctions with non-English prompts using various currencies.} 

% For each auction, we find evidence of LLM bids deviating relative to the risk-neutral equilibrium. Risk-aversion explains these deviations well but not completely. In particular, they don't explain over-bidding and under-bidding in the SPSB auction, where the equilibrium is in dominant strategies, so holds irrespective of risk attitudes. However, while bidding data appears to deviate from theoretical prediction dramatically, it does so in many of the same ways human bidding data deviates from theory.

% Our main results are plotted in Figure \ref{fig:fpsb-plan} and Figure \ref{fig:tpsb-plan}, alongside the prediction of rational theory. On each plot, we fit a Loess-smoothed curve -- these curves impose no parametric restrictions on the data and constitute evidence of monotonicity. A formal test is also offered later this section. We will first discuss the results of each of our simulations compared to the existing literature, highlighting five points in the analysis. 1) LLM play, in aggregate, looks similar to human play. 2) However, individual bids from LLMs are often mistakes, far from the optimal play. 3) These mistakes are consistent with risk aversion. This argument is made precise by identifying CARA and CRRA utility coefficients for LLM bidding and comparing them to the literature. 4) Mistakes in the SPSB, though unattributable to risk attitudes, also occur at the same rate as in human literature with some qualitative differences. 5) In sum, we find the largest qualitative difference between human and LLM bidding is that LLM bidding is higher variance. This motivates some interest in prompting which we return to later in the paper. 

\subsubsection{Simulation evidence}

For each of the IPV auction formats, we perform 5 experiments with 15 rounds and 3 LLM bidders competing over a prize.\footnote{While these specifications are designed to meet the `classic' auction formats of the literature, we also report the results of several variants in Appendix \ref{app:robustness}. In particular, we run one-shot auction simulations (as opposed to multi-round) with $n = 4, 5,$ and $10$ bidders, auctions with prompts using various currencies, and auctions with non-English prompts using various currencies.}

We see that the sealed-bid formats demonstrate evidence of monotone bidding, which is one of
the most stable hallmark of experimental lab results on auctions. This observation is tested with
Kendall’s \(\tau_b\) test for monotonicity, the results of which are reported in Table \ref{tab:risk-neutral-eqm}.

For the FPSB, the Loess-smoothed data curve is higher than the BNE prediction, matching closely
with the existing empirical evidence for first-price auctions \cite{cox1988theory}. For SPSB and TPSB auctions, the Loess-smoothed bidding curves appear generally lower than expected, relative to
both empirical and theoretical benchmarks. The lowest Loess-smoothed bidding line corresponds
to the all-pay auction.
Though bidding strategies in aggregate roughly correspond to empirical benchmarks (especially
when they diverge from theory), our results have high variance, suggesting a wide span of bidding
plans from individual LLMs. Uniformly low bids are also prevalent across the range of bidder values.

  

\begin{figure}[hbtp]
    \centering \includegraphics[width=\linewidth]{Figs/SPFP.png}
    \caption{\textbf{ Comparison of FPSB(Left)  and SPSB(Right) under IPV setting.} Bidders' values are independently drawn from a uniform distribution of [0, 99] and LLM bids given a value are represented by blue circles. The black dotted line represents the 45$^\circ$ line (the dominant strategy for SPSB auctions), the green dashed line represents the Bayes-Nash Equilibrium on the FPSB plot, and the solid red line represents the LOESS-smoothed bidding curve for both plots. }\label{fig:fpsb-plan} 
\end{figure}

\begin{figure}[hbtp]
    \centering \includegraphics[width=\linewidth]{Figs/TPAL.png}
    \caption{\textbf{ Comparison of TPSB and all-pay under IPV setting.} Bidders' values are independently drawn from a uniform distribution of [0, 99] and the experimental data points are represented by blue circles. The black dotted line represents  bid=value (the dominant strategy for SPSB auctions), the green line on the FPSB figure denotes the Bayes-Nash Equilibrium, and the solid red line represents the LOESS-smoothed bidding data. Here, we show results for third-price auctions on the left, and all-pay auctions on the right.
    \label{fig:tpsb-plan}}
\end{figure}

% First, we summarize the distributional features of our sealed-bid auctions. Bidding is diffuse, with most bidders deviating from optimal play. Table \ref{tab:bidding_errors} summarizes these distributions, categorizing bids as underbids, value bids, or overbids relative to one's value. However, aggregate patterns across auction formats provide encouraging evidence that, in aggregate, LLM agents are responding strategically to the change in auction format in ways theory predicts. First, we find evidence of monotonic bidding (in aggregate), testing this with a Kendall-$\tau$ test in Table \ref{tab:kendall_tau}. Monotonicity is an empirical regularity in the auctions literature, and the basis for the ubiquitous monotonicity assumption taken in auction theory. Second, we find evidence that as $k$ increases for the $k$-th price auction, the distribution of bids versus values shifts upward \citep{monderer2000k}. In the FPSB auction, only 4.4\% of bids exceed the bidder's value (a weakly dominated strategy), with the vast majority of bidders underbidding relative to their values. In contrast, the distribution shifts upward for SPSB (12.44\% bids higher than their value) and TPSB (27.11\% bids higher than their value). $\chi^2$ tests between the bidding distributions confirm this visual intuition, with results reported in Table \ref{tab:chi_square}. The closest evidence from the empirical literature, \citet{kagel1993independent}, find similar results, also running $\chi^2$ tests to confirm that distributions are distinct and moving relatively up from the FPSB to the SPSB to the TPSB auction. The difference between \cite{kagel1993independent}'s results and our own on this matter is one of magnitudes -- while the bidding distributions for the $k-$price auction are significantly distinct and increasing in $k$, there is more separation between bidding distributions with human bidders than LLM bidders.

\begin{table}[h]
\centering
\renewcommand{\arraystretch}{1.4}

\begin{tabular}{lccc}
\hline
\textbf{Auction Type} & \textbf{Underbidding (\%)} & \textbf{Value Bidding (\%)} & \textbf{Overbidding (\%)} \\
\hline
First-Price Sealed-Bid  & 82.22  & 13.33  & 4.44  \\
Second-Price Sealed-Bid  & 65.78  & 21.78  & 12.44 \\
Third-Price Sealed-Bid  & 42.22  & 30.67  & 27.11 \\
All-Pay & 94.67  & 4.44   & 0.89  \\
\hline
\end{tabular}
\caption{\textbf{Bidding relative to one's value.} Underbidding is defined as $\beta_i(v) < v_i$; Value bidding is defined as $\beta_i(v) = v_i$; Overbidding is defined as $\beta_i(v) > v_i$.\label{tab:bidding_errors}}
\end{table}

\begin{table}[h]
    \centering
    \begin{minipage}{0.48\textwidth}
        \centering
        \begin{tabular}{llcc}
            \toprule
            Auction 1 & Auction 2 & $\chi^2$ Statistic & p-value \\
            \midrule
            First-Price  & Second-Price  & 22.97  & 0.0063 \\
            First-Price  & Third-Price   & 35.68  & 4.52e-05 \\
            Second-Price & Third-Price   & 36.67  & 3.02e-05 \\
            All-Pay      & First-Price   & 11.77  & 0.2642 \\
            All-Pay      & Second-Price  & 24.69  & 0.0033 \\
            All-Pay      & Third-Price   & 59.46  & 1.70e-09 \\
            \bottomrule
        \end{tabular}
        \caption{\textbf{Chi-Square Test Results for Auction Comparisons.} This table presents the chi-squared test statistics and p-values comparing different auction formats.}
        \label{tab:chi_square}
    \end{minipage}
    \hfill
    \begin{minipage}{0.42\textwidth}
        \centering
        \begin{tabular}{lcc}
            \toprule
            Auction Type & Kendall's $\tau_b$ & p-value \\
            \midrule
            First-Price  & 0.672*** & $<$ 0.0001 \\
            Second-Price & 0.666*** & $<$ 0.0001 \\
            Third-Price  & 0.671*** & $<$ 0.0001 \\
            All-Pay      & 0.660*** & $<$ 0.0001 \\
            \bottomrule
        \end{tabular}
        \caption{\textbf{Kendall's $\tau_b$ test for bid monotonicity.} $\tau_b = 1$ indicates perfect positive monotonicity, while $\tau_b = -1$ indicates perfect negative monotonicity.}
        \label{tab:kendall_tau}
    \end{minipage}
\end{table}

% Despite features in aggregate corresponding to rational theory, most individual bidding is not optimal play. We can account for many of these mistakes by modeling risk preferences. To quantify risk attitudes in LLM bidding, we estimate risk coefficients under CARA and CRRA preferences for the FPSB, TPSB and all-pay auctions according to the specification in Section \ref{sec:ipv-theory}. The SPSB is excluded from this analysis as dominant strategy equilibria are invariant to risk attitudes. Results are summarized below in Table \ref{tab:risk_aversion}.

% \begin{table}[h]
% \centering
% \renewcommand{\arraystretch}{1.4}
% \setlength{\tabcolsep}{10pt}

% \begin{tabular}{lcccc}
% \hline
% \textbf{Model} & \textbf{FPSB} & \textbf{TPSB} & \textbf{All-Pay} & \textbf{Literature Range} \\
% \hline
% CRRA ($\rho$)  & 0.49  & 1.01  & 0.52  & 0.3–0.8  \\
% CARA ($\lambda$)  & 0.011  & \avs{}  & \avs{}  & 0.008–0.02  \\
% \hline
% \end{tabular}

% \vspace{0.3cm}
% \caption{\textbf{Estimated risk aversion parameters.} Literature range is based on \citet{cox1988theory} for FPSB and \citet{dechenaux2015survey} for All-Pay auctions. For TPSB, we compare to the limited experimental evidence from \citet{kagel1993independent}.}
% \label{tab:risk_aversion}
% \end{table}

% The most robust empirical finding from the current literature is for the FPSB auction; it's a well-known empirical result that bidders in the FPSB bid above the risk-neutral Nash equilibrium level due to risk aversion, thereby failing revenue equivalence \citep{cox1988theory, kagel1993independent, kagel2020handbook}.
% The implied CARA and CRRA parameters align closely with measured risk attitudes in human subjects. 
% For TPSB auctions, estimated risk-aversion parameters are notably larger than both the FPSB parameters and typical literature benchmarks, but still consistent with a coherent risk-averse framework. One plausible explanation for the higher observed risk-aversion in TPSB may be the compounding influence of loss aversion, which biases bids in the same downward direction as risk aversion. The all-pay auction results present a more complicated picture. The estimated parameters for all-pay auctions suggest significantly higher risk aversion. However, these estimates cannot fully explain observed deviations: low-valued bidders, for instance, systematically fail to reduce their bids below equilibrium predictions, contrary to standard risk-aversion reasoning. This inconsistency may indicate a deeper misunderstanding among LLM agents about the strategic complexity inherent to the all-pay mechanism. It is worth noting that human experiments on all-pay auctions also yield highly variable results, reflecting the cognitive challenge this format presents to both human and LLM bidders alike \citep{gneezy2006all, dechenaux2015survey}.

 As with human bidders, LLM bidding in the SPSB auction exhibits numerous departures from equilibrium play. Table \ref{tab:bidding_errors} summarizes the frequency of underbidding, value bidding, and overbidding by LLM bidders, comparing these results directly with human data from \citet{kagel1993independent}. In our experiments, $21.78\%$ of bids exactly match bidders' values, closely aligning with the $27\%$ reported by \citet{kagel1993independent}.\footnote{Note this comparison involves slight measurement differences -- \citet{kagel1993independent} allow a small margin (within $\$0.05$) for classifying bids as ``truthful," whereas we require exact equality.} Despite similar overall rates of mistakes, the nature of these deviations differs strikingly. In a SPSB auction with $5$ players, \citet{kagel1993independent} find the distribution of bids to be right-skewed, with most of the bidding mass (67\%) on over-bidding one's value. In contrast, we find the distribution of bids to be primarily left-skewed, with most of the bidding mass (66\%) on under-bidding one's value. The contrast highlights that despite both under-bidding and over-bidding being violations of strategy-proofness, they represent two distinct types of behavioral errors empirically -- an insight that has not been well rationalized by the auctions literature yet. The prevalence of underbidding among LLM bidders further suggests loss-aversion to a degree not observed in human bidders. And, taken as a whole, we see the mistakes under SPSB as additional motivation for exploring the effect of obviously strategy-proof designs.

In this section, we've highlighted several features of individual and aggregate LLM bidding behavior. Perhaps the feature of LLM data that makes it the hardest to interpret individual bids is the increased variance. We suggest this increased variance in LLM bidding is due to increased random play by LLMs, due both to a desire to gather exploratory data as well as to disrupt opponents' predictions. The ability to peer into LLMs' chain-of-thought reasoning makes this particularly transparent. For instance, LLM bidders frequently articulate plans like the following from Bidder betty in a round of the FPSB auction: 


\begin{quote}
    "I'll adopt a balanced strategy, bidding $65\%$ of my value. \textbf{I'll also introduce random bids occasionally to disrupt predictability}. Monitoring competitor's bids remains essential to adjust my strategy accordingly."
\end{quote}


This ability to randomize is just one way in which LLMs may differ from human bidders, but there are many others -- say, the ability to reliably calculate 65\% of one's value or to monitor and perform inference over the full history of competitor bids. The variance in LLM play suggests to us that looking at LLM data in aggregate, rather than scrutinizing individual data points, may be more productive. Just as individual human bidders encode their own biases and tendencies into their actions, no individual LLM bidder's behavior is particularly instructive -- however, in aggregate, auction data with LLM data encodes behavioral patterns that may align closely with many of the empirical regularities from auction theory. It is our view that documenting these points of alignment or divergence will make the future use of LLM data for auction theory or mechanism design more productive.

Finally, one may ask how prompting may change the general direction of our results. This is a first-order concern for the eventual use of synthetic data in mechanism design and experimentation. To address this, we pursue two additional lines of inquiry:
%
\begin{itemize}
    \item[(1)] We test prompt robustness. We implement our FPSB and SPSB environments with prompts in a variety of prompt styles, in different languages and with different currencies to study how bids change under prompt variation. These results are reported in Appendix~\ref{app:prompt}. We find that naive changes in prompt style, language or currency do not qualitatively change any of the results reported in this section.
    %
    \item[(2)] Beyond naive prompt changes, we test prompts which influence the LLM agents' understanding of the underlying economics. To do this, we design prompts to teach LLMs particular mental models or avoid particular failure modes; e.g., a prompt to consider Nash deviations, to conceptualize the auction as an extensive-form game, or a prompt which encourages agents to bid in a risk-neutral manner. These results are reported in Section~\ref{session:Intervention}. We find that \textit{these} interventions can dramatically change our results. 
\end{itemize}

Our results suggest to us that prompting \textit{can} induce dramatically different behavior in LLMs by making the underlying logic of mechanisms more apparent, but is unlikely to generate different results with naive changes in prompts. It remains an open engineering question as to `how much' prompting is required to match distributions of existing empirical evidence more closely.  

% Second, as with human bidders, LLM bidding in the SPSB auction exhibits numerous deviations from the dominant strategy of truthful bidding. In our bidding data, 22\% bids in the SPSB auction are at their value, and in \citet{kagel1993independent}, 27\% are.\footnote{Note this is a fuzzy comparison -- \citet{kagel1993independent} introduce some fault tolerance, with any bid within \$0.05 cents of a value being tagged as bidding one's value. We tag bids as `value bidding' only if agents bid \textit{exactly} their value.} However, the shape of these deviations differs interestingly. In a SPSB auction with $5$ players, \citet{kagel1993independent} find the distribution of bids to be right-skewed, with most of the bidding mass (67\%) on over-bidding one's value. In contrast, we find the distribution of bids to be primarily left-skewed, with most of the bidding mass (66\%) on under-bidding one's value. The contrast highlights that despite both under-bidding and over-bidding being violations of strategy-proofness, they represent two distinct types of behavioral errors empirically -- an insight that has not been well rationalized by the auctions literature yet. Table \ref{tab:bidding_errors} summarizes these results below.

% Of all the classic sealed-bid auctions, the literature on the all-pay auction has the most variance in documented experimental behavior ~\citep{dechenaux2015survey}. However, as in the FPSB auction, one major feature of the empirical literature is overbidding relative to the BNE, while still bidding below one's value (overbidding one's value is a dominated strategy). This holds under a variety of auction specifications in the literature, and is echoed in our results: $89.78\%$ of bids are greater than the implied BNE bid but less than the agent's value, and in particular, there is almost no overbidding one's value (with $0.89\%$ of agent's bids greater than their value). Risk-aversion is the literature's preferred explanation for overbidding relative to the BNE. However, some experiments in the literature also report substantial heterogeneity in bidding strategies, with bidding strategies differing sharply based on bidder values or types. Most commonly this is reported as a bimodal bidding distribution, with low-valued bidders bidding low and high-valued bidders bidding high relative to the BNE (e.g., \citet{gneezy2006all}). Our results consistently dominate the BNE prediction (in a way that corresponds to the general results of the literature) but demonstrate no heterogeneity between bidders with high or low values.

% In fact, with the exception of the SPSB auction, risk-aversion provides a compelling explanation for LLM bidders' deviations from risk-neutral theory -- LLM bidders demonstrate risk-aversion, and because of that fail to bid according to each game's BNE. To make this more precise, we fit CARA and CRRA utility functions to quantify this risk-aversion and compare our findings to the literature, summarized in Table \ref{tab:risk_aversion}. In particular, experimental evidence has consistently shown that bids are shaded upward towards one's value in the FPSB auction relative to risk-neutral predictions \citep{cox1988theory}, a finding our results replicate. Agents are willing to give up some surplus by bidding closer to their true value in exchange for a higher probability of winning the auction. Our estimated CRRA coefficient of $\rho_{FPSB} = 0.42$ falls within the range of 0.33 to 0.67 commonly reported in the experimental literature, suggesting LLM agents exhibit risk-aversion at levels comparable to human subjects. Similarly, risk-aversion provides compelling results for the TPSB. Theory predicts that risk-averse agents will shade bids downwards towards their value. The lone empirical study on the TPSB corroborates this prediction and aligns with our results, wherein LLM agents underbid the risk-neutral BNE prediction by bidding towards their value. Our estimated CRRA coefficient of $\rho_{TPSB} = 1.01$ suggests agents are more risk-averse on average in the TPSB auction than in the FPSB auction, with high variance in the bidding distribution.

% However, while LLMs demonstrate risk-aversion in a way which seems similar to the empirical literature, the data with LLM bidders is also qualitatively different from that with human bidders. Most of the difference seems attributable to the higher variance in their bidding behavior. We suggest this high variance stems from increased random play by LLMs due to both data gathering and attempts to disrupt opponents' strategies. The ability to peer into LLM's chain-of-thought reasoning makes this explanation legible. For example, LLMs frequently make plans of the sort exhibited by bidder Betty in the first run of the FPSB auction:

% \begin{quote}
%     "I'll adopt a balanced strategy, bidding 65\% of my value. I'll also introduce random bids occasionally to disrupt predictability. Monitoring competitor's bids remains essential to adjust my strategy accordingly."
% \end{quote}

% This emergent strategic randomization suggests that looking at the shape of this data in aggregate, rather than scrutinizing individual data points, provides more meaningful insights. Just as individual human bidders encode their own biases and tendencies, no individual LLM bidder's behavior is particularly instructive in isolation. However, in aggregate, the auction data encodes behavioral patterns that align with established auction theory.

% Finally, one may ask how much prompting may change the general direction of our results. This is a first-order concern for the eventual use of synthetic data in mechanism design and experimentation more broadly. To address this,
% we pursue two additional lines of inquiry:
% %
% \begin{itemize}
%     \item[(1)] We test prompt robustness. We implement our FPSB and SPSB environments with prompts in a variety of prompt styles, in different languages and with different currencies to study how bids change under prompt variation. These results are reported in Appendix~\ref{app:prompt}. We find that naive changes in prompt style, language or currency do not qualitatively change any of the results reported in this Section.
%     %
%     \item[(2)] Beyond naive prompt changes, we test prompts which influence the LLM agents' understanding of the underlying economics. To do this, we design prompts to teach LLMs particular mental models or avoid particular failure modes; e.g., a prompt to consider Nash deviations, to conceptualize the auction as an extensive-form game, or a prompt which encourages agents to bid in a risk-neutral manner. These results are reported in Section~\ref{session:Intervention}. We find that \textit{these} interventions can dramatically change our results. 
% \end{itemize}

% Our results suggest to us that prompting \textit{can} induce dramatically different behavior in LLMs by making the underlying logic of mechanisms more apparent, but is unlikely to generate different results with naive changes in prompts. It remains an open engineering question as to `how much' prompting is required to match distributions of existing empirical evidence more closely.  

% Finally, there are two additional questions that we find interesting but do not run further analysis on. One, while we 

% \hline\hline \avs{I want to say the stuff at the bottom. it seems important. this section also seems a bit like word salad -- the best way to summarize? a table at the end? organize the arguments of the section as a whole?}

% We also find it interesting that, looking in aggregate, there is also a general tendency to underbid one's value not quite seen to the same extent in the literature with human bidders. This underbidding tendency could potentially be explained by greater loss aversion in LLMs relative to humans, though further investigation would be needed to confirm this hypothesis.

% The Loess-smoothed curves demonstrate evidence of monotone bidding. This is statistically tested with Kendall's $\tau_b$ test for monotonicity, the results of which are reported in Table \ref{tab:kendall_tau}.
% \begin{table}
%     \centering
%     \begin{tabular}{lcc}
%         \hline
%         Auction Type & Kendall's $\tau_b$ & p-value \\
%         \hline
%         IPV First-Price  & 0.672*** & < 0.0001 \\
%         IPV Second-Price & 0.666*** & < 0.0001 \\
%         IPV Third-Price  & 0.671*** & < 0.0001 \\
%         IPV All-Pay      & 0.660*** & < 0.0001 \\
%         \hline
%     \end{tabular}
%     \vspace{.3 cm}
%     \caption{\textbf{Kendall's $\tau_b$ test for bid monotonicity in the IPV setting.} Here,
%     $\tau_b = -1$ corresponds to perfect negative monotonicity; $\tau_b = 0$ corresponds to no statistically observable evidence of monotonicity; and $\tau_b = 1$ corresponds to perfect positive monotonicity. A value of $.67$ means that 67\% of the time, when comparing two random bidders, the one with a higher private value also has a higher bid.
%     \label{tab:kendall_tau}}
% \end{table}

% The results in Table \ref{tab:kendall_tau} provide strong statistical evidence for monotonic bidding across all auction formats. For each format, approximately two-thirds of pairwise comparisons exhibit the expected positive relationship between values and bids, consistent with the monotonicity assumption in auction theory and empirical findings with human bidders.
% To further quantify the extent of risk aversion in our FPSB results, we fit a linear model relating bids to values and estimate the CRRA coefficient, presented in Table \ref{tab:crra_results}.
% % \begin{table}[h]
% % \centering
% % \caption{Constant Relative Risk Aversion (CRRA) Estimation Results}
% % \label{tab:crra_results}
% % \begin{tabular}{lc}
% % \hline
% % \multicolumn{2}{c}{Dependent Variable: Bid} \
% % \hline
% % Constant & $0.8369$ \
% % & $(0.8656)$ \
% % Value     & $0.7031^{}$ \
% % & $(0.0000)$ \
% % \hline
% % Observations & 222 \
% % R-squared    & 0.6566 \
% % \hline
% % \multicolumn{2}{l}{CRRA Coefficient Estimates:} \
% % Mean of Individual Estimates & $0.4205$ \
% % Median of Individual Estimates & $0.5395$ \
% % Risk-Neutral Reference Point & $0.6667$ \
% % \hline
% % \multicolumn{2}{l}{Notes: $p$-values in parentheses.} \
% % \multicolumn{2}{l}{$^{}$Significant at the 1% level.} \
% % \end{tabular}
% % \end{table}
% The estimated value coefficient of 0.7031 is higher than the risk-neutral prediction of 0.6667, yielding a CRRA coefficient of approximately 0.42. This confirms that LLM agents exhibit risk-aversion in the FPSB auction, bidding more aggressively than predicted by risk-neutral theory but at levels consistent with human behavior in laboratory experiments.
% In summary, our LLM agents replicate many of the key empirical regularities documented in the auction literature: monotonic bidding strategies, systematic deviations from risk-neutral predictions consistent with risk-aversion, and strategic responsiveness to changes in auction format. The major exception is the pattern of errors in the SPSB auction, where our LLM agents tend to underbid rather than overbid as commonly observed with human subjects. This difference merits further investigation and may reflect fundamental differences between human and LLM decision-making processes in strategy-proof environments.


\subsection{Obviously Strategy-proofness}\label{session:OSP}
For the rest of the paper, we restrict our attention to the second-price auction mechanism. Having considered LLM play in a variety of sealed-bid auctions, we now vary the representation of the second-price auction to make the auction `easier' to play. Namely, we turn our attention to clock auctions. Clock auctions are dynamic auction formats where the price increases incrementally (or decreases incrementally in reverse auctions), and bidders decide whether to stay in or drop out at each price point.
% These auctions are strategy-proof mechanisms since a player does not need to form beliefs about other players' values and types; 
% she can simply report her value truthfully (bid her value) because she can do no better regardless of what the other players bid or what their bidding strategies are. While truthful reporting is a weakly dominant strategy here, many people in reality do not report their types truthfully in strategy-proof mechanisms. 
\citet{li2017obviously} provides an argument as to why clock auctions may be easier to recognize as strategy-proof for humans, and formalizes this as an equilibrium refinement on strategy-proofness called {\em obvious strategy-proofness} (OSP).
% A mechanism is OSP if it is evident to a bidder that truthful bidding is optimal, regardless of others' actions, at every decision point.
% OSP mechanisms are simple to play and thus have a lot of desired properties. 


\subsubsection{Setting}
There are three bidders in each auction but now, following \citet{li2017obviously}, bidders draw {\em affiliated private values} (APV) of the form $v_i = c + p_i$. The common component $c$ is drawn once, $c \sim U[0, 20]$, and the private component $p_i$ is drawn $n$ times independently, $p_i \sim U[0, 20]$. Winners of the auction receive their own value of the prize $v_i$ when they win, and the `common' and `private' components  serve to make values correlated. We switch to considering APV settings in this section for two reasons. First, in spirit of our larger goal to provide many benchmarks of LLM play in various environments. Second, to enable comparison with the existing empirical literature; the two largest experiments of which have been conducted in APV settings. 

The {\em ascending clock auction} (AC) has
a global price clock that increases in some fixed increment and bidders who choose when to drop out. The last bidder remaining wins the prize at the price where the second to last bidder dropped out. If multiple bidders are the last to drop out, we resolve ties by randomly allocating the prize among them and setting the winning price to that mutual drop out price. We also study the {\em blind ascending clock auction} (AC-B), a variant 
where bidders are not told when  other bidders   drop out.
%The SPSB auction is as defined in Section \ref{section:classic}. 

\subsubsection{Theoretical benchmarks}
What effect an affiliated values environment has on strategic behavior may not be ex-ante obvious. On one hand, the correlation in values may strengthen incentives for extremal-valued bidders, so that extremal-valued bidders are actually \textit{more} likely to bid their value even in the sealed-bid format (echoing arguments from \citet{li2017obviously}). On the other hand, APV may be cognitively more complex for bidders relative to IPV environments, thereby inducing errors in play. In either case, the three auction formats in this section are strategically equivalent, and so the affiliation is, in the rational sense, a red herring: for all three auctions it is still a dominant strategy to bid (or drop out at) one's value. The SPSB, AC-B and AC auctions are strategically equivalent and share an equilibrium in dominant strategies:
\[{\beta}^{**}(v) = v.\]

An agent's bid, in the ascending clock format, is equivalent to the price at which they drop out of the auction at. We will use the terms `bid' and `drop-out price' interchangeably in the context of clock auctions.
The equilibrium refinement of OSP, introduced by \citet{li2017obviously}, provides a formal framework for analyzing the cognitive simplicity of mechanisms.
Consider a strategy $S_i$ for player $i$. 
We say that $S_i$ {\em obviously dominates} another strategy $S'_i$ if, for any deviation from $S_i$ to $S'_i$, the best possible outcome from $S'_i$ is no better than the worst possible outcome from $S_i$.

Formally, for all histories $h$ where both $S_i$ and $S'_i$ are consistent with $h$:
\begin{equation*}
\min_{z \in Z(S_i, h)} u_i(z) \geq \max_{z' \in Z(S'_i, h)} u_i(z')
\end{equation*}
where $Z(S_i, h)$ is the set of terminal histories that can result from playing $S_i$ after history $h$ and
$u_i(z)$ is player $i$'s utility for terminal history $z$.  
% \dcp{explain in words what a `history` is in a clock auction} 
In a clock auction, the history refers to the sequence of events that track bidder activity throughout the auction. Specifically, it records whether and when bidders drop out as the price increases, providing insight into bidding behavior and competition dynamics. A mechanism is said to be obviously strategy-proof if it has an equilibrium in dominant strategies.  Both ascending clock auctions studied in this paper are obviously strategy-proof, though the AC-B auction provides bidders with less information than the AC auction by suppressing when bidders drop out.


\subsubsection{Empirical benchmarks and results}
For each of the APV auction formats, we perform 5 experiments with 15 rounds and 3 LLM bidders competing over a prize. A single round for the clock format is the run of a global clock from the initial price of $\$0$ to the final clearing price. The complete simulation of a single clock round is provided in \ref{app:dynamic}.

Our main results are plotted in Figure \ref{fig:osp_bid}. We find that the clock formats induce play significantly closer to one's value than the sealed-bid format. Table \ref{tab:auction-comparison} summarizes this concisely: the clock formats, aggregated, have an average difference between bid and value of $1.83$ (with $73.9\%$ of agents dropping out at their value, excluding winners), while the SPSB bidding has an average difference between bid and value of $6.26$ (with $18.7\%$ of all bids equal to an agent's value). This corresponds closely with \citet{li2017obviously}'s experimental results: aggregated across all rounds, he reports a mean difference between bid and value of $2.48$ for the ascending clock auction, and a mean difference of $6.52$ for the sealed-bid auction. In general, the clock format improves play for LLM bidders in much the same way as with human bidders.

\begin{figure}[h]  
    \centering  
\includegraphics[width=\linewidth]{Figs/OSP_gpt4.png}  
    \caption{\textbf{
    Comparison of three strategically equivalent auctions (the AC, AC-B and SPSB auctions) in an APV setting.} 
    The AC and AC-B auctions are obviously strategy-proof, whereas SPSB auctions are only strategy-proof. The black dot-dash line represents the dominant strategy where a bidder's bid equals their value. In the clock auctions, red triangles indicate non-winners' drop out prices and blue circles indicate the final price that winners win at. In the SPSB auction, red triangles indicate non-winners' bids and the winning bid in clock auctions. Note: the blue dots are not expected to equal an agents value in the clock auctions, as we do not observe a winner's drop-out price in the clock auction. Bidders drop out close to their value in clock auctions, while bidders bid quite diffusely in the corresponding sealed-bid auction.
    \label{fig:osp_bid}}
\end{figure}

\begin{center}
    \begin{table}[h]
\centering
\begin{tabular}{l|c|c|c}
\hline
\textbf{Metric} & \textbf{AC} & \textbf{AC-B} & \textbf{SPSB} \\
\hline
Truthful Bidding (\%) & 70.0\% & 76.7\% & 18.7\% \\
\hline
Average Difference & 2.03 (0.39) & 1.63 (0.40) & 6.26 (0.62) \\
\hline
Sample size& 90 & 90 & 150 \\
\hline
\end{tabular}
\vspace{.5 cm}
    \caption{\textbf{Share of truthful bidding in the AC, AC-B and SPSB auctions in an APV setting}.  
    Following the definition in \citet{li2017obviously}, we define truthful bidding in the ascending clock (AC) and ascending clock with bid (AC-B) auctions as dropping out within one increment of the bidder's true value (excluding winners), and in the second-price sealed-bid (SPSB) auction as submitting a bid exactly equal to one's true value. The distinction reflects the strategic differences inherent in each mechanism: in clock auctions, bidding is a dynamic process where drop-out behavior approximates valuation, while in SPSB auctions, truth-telling is a dominant strategy only when the bid precisely equals the true value. We find that the SPSB auction exhibits a substantially lower rate of truth-telling under this definition.}
\label{tab:auction-comparison}
\end{table}
\end{center}

\vspace{-.5cm}

\citet{breitmoser2022obviousness} provide additional evidence that clock auctions improve truthful play, and also provide evidence that play under AC is slightly better than under AC-B. While we replicate the former result, we fail to replicate the latter result. Play under the AC auction was not significantly different compared to the AC-B auction ($t = 0.71, p =0.48$) in our experiments. We report all the two-sample t-tests of the mean absolute difference between bids and values in Table \ref{tab:t-test-results} -- both AC and AC-B mechanisms show significantly smaller bid-value deviations than the SPSB auction (p $<$ 0.001), but not from one another. 

\begin{table}[h]
\centering
\begin{tabular}{l|c|c}
\hline
\textbf{Comparison} & \textbf{t-statistic} & \textbf{p-value} \\
\hline
AC v.s. AC-B & 0.71 & 0.48 \\
\hline
AC v.s. SPSB & -6.28 & 3.21e-09 \\
\hline
AC-B v.s. SPSB& -5.78 & 4.09e-08 \\
\hline
\end{tabular}
\caption{\textbf{Two-sample t-tests comparing the mean absolute difference between bids and values in the AC, AC-B and SPSB auctions in an APV setting}. Results demonstrate that both clock auctions have lower mean absolute difference between bid and value vs the sealed-bid auction. However, the clock auctions themselves are not significantly different. Bid data for clock auctions is excluding winners.
\label{tab:t-test-results}}
\end{table}

Finally, we observe that the clock format improves on a particular failure mode of the SPSB—namely, bidding \textit{above} one's value. The result is that the share of total mistakes in clock auctions overwhelmingly comes from dropping out too early rather than too late. Although OSP mechanisms cannot rationalize a difference between over- and under-bidding (both represent deviations from an obviously dominant strategy), this evidence further suggests that over- and under-bidding represent categorically \textit{distinct} failure modes. Further exploration into rationalizations of the difference between these two failure modes remains an interesting and open question for future research.


% \dcp{suprising that the SPSB results so different from IPV SPSB results. worth a discussion and begs
% the question why you didn't also do an IPV clock treatment.}

\section{Ebay Auction Design}\label{session:ebay}

Having tested LLMs against classic results in auction theory, we  build an auction environment to test components of auction design in a setting inspired by eBay's online marketplace. eBay's auction platform, studied by \citet{einav2011learning} and others, features various design elements -- for example, a hidden reserve price (as opposed to a reserve price)
and proxy bidding -- that are \emph{in theory} less critical to auction outcomes
%(often being revenue-equivalent) given they induce bidders to bid their value,
yet \emph{in practice} appear to have substantial effects on bidder behavior. These differences highlight the gap between textbook theory and real-world bidding, where seemingly small design details can alter bidder behavior and thus price discovery and revenue.

% Having documented encouraging results from surveying LLM agents in classic auction designs, we next explore how LLM agents respond when confronted with more realistic conditions inspired by eBay's online marketplace. 

In this section, we simulate a stylized version of the eBay auction design to consider two design parameters more closely: the hidden reserve price and the implementation of a `closing rule'. The core mechanics of an eBay auction are simple: bidders on eBay use a proxy bidding system, which means that each bidder submits their ``maximum bid" to the platform, and the platform automatically bids on their behalf, only as much as necessary to keep the winning bidder in the lead (i.e., the second-highest bid + \$1, with incremental bidding managed automatically).
Bidders can increase their maximum bid at any point in time, until some closing time of 
the auction, giving this the flavor
of an incremental second-price auction.

% As such, if bidders play in the way eBay expects them to, the hidden reserve and the closing rule shouldn't change bidder strategy -- all bidders report their value as soon as they see the auction, and the auction does not update any further. 

The use of reserve prices in auctions is well understood in auction theory. Notably, 
\citet{myerson1981optimal} showed that setting an optimal reserve price can maximize expected revenue in certain
private value settings (for ``regular" distributions).
However, eBay also allows sellers (for a fee) to set a {\em hidden} reserve price: a threshold that bidders do not know until it is met by the highest maximum bid. 
%
While there is some work suggesting a hidden reserve can increase expected revenue in common value settings (e.g., \citet{vincent1995bidding}), in a private values environment with rational bidders using proxy bidding, the hidden reserve should not change behavior relative to a known reserve or even no reserve: bidders should still submit their maximum value upfront. While it is impossible to say whether bidders are bidding their true values, a large empirical study of hidden reserves on eBay suggests that bidders do not change bidding behavior much in the face of hidden vs.~non-hidden reserves \citep{einav2011learning}. 

A more dramatic departure from simple, truthful bidding is `bid sniping,' or the practice of placing bids in the final moments of an auction to avoid triggering competitive responses.
Sniping is prevalent on eBay, with dedicated e-snipe services automating last-second bid submissions.\footnote{E.g., \url{https://www.esnipe.com/}.} This phenomenon has been studied for decades, with explanations ranging from best responses to irrational types to behavioral inertia \citep{roth2002last}. Such wide prevalence of sniping may be especially troubling to economists, as it inhibits a core function of auctions: their ability to induce timely, efficient price discovery. One avenue to curb sniping may be a `closing rule' (also called a `soft-close' or `auction extension' rule) -- namely, when a new maximum bid comes in at the last minute, extend the auction close to allow for more bidding. This design detail has been credited with the substantial difference in bidding dynamics between earlier auctions on Amazon.com, which had this auction-extension rule, and auctions on eBay, which do not. Specifically, \citet{roth2002last} and \citet{ockenfels2006late} show that much less sniping occured on Amazon.com, with bids on Amazon.com occuring throughout the auction compared to the empirical evidence of sniping on eBay. 

\subsection{Setting}

In each eBay simulation, there are three bidders. Without the modified closing rule, 
the auction runs for 10 bidding periods. Bidders draw private values uniformly and independently from $[0, 99]$. In regard to  the auction environment, we consider four treatments, varying the hidden reserve and the closing rule:
%
\begin{itemize}
    \item[T1]: Standard eBay auction with proxy bidding.
    \item[T2]: eBay auction with a modified (auction extension) closing rule.
    \item[T3]: Standard eBay auction with a hidden reserve price.
    \item[T4]: eBay with a modified (auction extension) closing rule and a hidden reserve price.
\end{itemize}

For each hidden reserve price design, we run
experiments with hidden reserves at $r = 40, 50$, and $60$.

We model continuous time in the synthetic, modified eBay auction environment
%dcp , as querying LLMs is, by necessity, in discrete time. We deal with this 
as follows: each period, bidders bid (i.e., decide whether to increase their maximum bid or  hold) in a common knowledge, randomly-permuted order (e.g., today Charles, then Alice, then Betty will bid; tomorrow, Alice, then Betty, then Charles). Moreover, each bidder can decide whether to increase their maximum bid with knowledge of the updated state of the auction based on the actions of bidders that precede them that period.
In the final period, 
bidding is simultaneous and
no bidder will know whether they are the last bidder (simulating that, in a continuous time auction, you cannot know if another bidder will get a bid in between you and the auction's close).
Under the modified closing rule, if a higher bid is placed in the final period such that there is a new winner, 
the auction is extended by an additional period. This process continues until no further bids are made.

With this stylized eBay framework we abstract away from additional eBay design features 
such as shipping fees and Buy-It-Now. There was no prompting of LLMs to induce sophisticated play as in Section \ref{session:Intervention}. The complete simulation process is explained in Appendix \ref{app:ebay} and the prompts are listed in Appendix~\ref{app:prompt}.

% The eBay marketplace also offers the Buy-It-Now (BIN) option, a posted price at which a buyer can short-circuit the auction by making an immediate purchase. 
% In a standard private-value auction model with symmetric bidders, sellers would not want to offer such an option, because a lower-valuation bidder could preempt a higher-valuation bidder, thereby reducing the expected price.

% Nonetheless, \citet{einav2011learning} find that the probability of a sale via the BIN option is low but non-zero, and it declines quickly at higher BIN prices. 
% Moreover, a higher BIN price is correlated with a higher final sale price, regardless of whether the item sells via the auction or the BIN option. 
% Existing literature suggests a BIN option can attract bidders who value immediate purchase or prompt early bidding that escalates.

% To study the effect of the short-circuit option, we design our experiment starting from a very low starting price and examining the role of this BIN feature among LLM participants.

% % \noindent
% \begin{quote}
% \textbf{BIN option prompt:}
% \textit{  This item may also has a "Buy It Now" (BIN) option. If available, you can purchase the item immediately for \{\{BIN\_price\}\}, ending the auction. }
% \end{quote}


% \subsection{Shipping Fees}
% We also compare the effect of a higher starting price against higher ex-post shipping fees. 
% Although classic auction theory suggests the two are equivalent, \cite{einav2011learning} find that increasing shipping fees reduces the sale price, but by less than a one-for-one ratio.

% % \noindent
% \begin{quote}
% \textbf{Shipping Fees Prompt:}
% \textit{  Shipping costs are \{\{shipping\_fees\}\}. This amount will be added to the final price if you win the auction. }
% \end{quote}


\subsection{Simulation Process of eBay Auction}
For eBay auctions, we discretize the continuous bidding period into 10 periods. 
Each period, LLM
agents decide whether to increase their bid or hold their current bid. 
The bidding process follows
a structured format, where bidders act in a predefined sequence each period (e.g., Charles, then
Alice, then Betty). 
For each round,
bidders are reminded of the auction rules, similar to those outlined in Section \ref{sec:process}.
Agents are informed of past price changes through a transcript, such as ``On day
1, the price changed to 1. On day 2, the price changed to 3.''
The detailed
prompt and an example of one clock are listed in \ref{app:ebay}.


\subsection{Simulation evidence}
\textbf{Hidden reserve price results.}

Figure $\ref{fig:ebay_revenue}$  shows that the auction revenue remains basically unchanged regardless of the experiment type. In Table $\ref{tab:revenue_tests}$, we report a table of two-sided t-tests making this explicit.  
% \dcp{discuss: I think this means your reserve is too low.}
% \dcp{strange. you'd expect reserves, whether hidden or not, to increase revenue, and especially with just 3 bidders. you should comment on this, and try higher reserves. you can use Myerson's formula to calculate the optimal reserve.} 
%

\noindent\textbf{Closing rule results.}
Second, we demonstrate that modifying eBay's closing rule to allow for an  extension of the auction when there is final-period bidding  
has a major effect on bidder behavior. 
We define the final-winning-bid time as the last period in which the eventual winner changes her maximum bid. If she raises her maximum in a later period, that later period becomes the recorded time.
% For this, we say the `final winning bid time' is the period at
% which the winner is settled; i.e., the period in which the maximum bid was placed by the winning bidder. \dcp{what if the winning bidder increased their maximum bid in some later period? would this be the final winning bid time, or would it be tje period in which they first bid enough to win? be precise.}
%
This is a measure for when the bidder who would win the auction was determined.
% \dcp{`decided` is not quite clear.}
%
Figure $\ref{fig:final_winning_bid_shift}$ reports the `final winning bid time' as a histogram across auction instances with the modified closing rule and without the modified closing rule. Figure $\ref{fig:ebay_cdf_final_winning_bid}$ reports a CDF across the four treatment types (T1 - T4). 
Without the modified closing rule, we see that LLMs predominantly bid in the period before close; however, with the modified closing rule, this behavior disappears. 
In Figure~$\ref{fig:ebay_cdf_final_winning_bid}$, we see the same behavior. 
With the modified closing rule, the proportion of last-period bids drops substantially; adding a hidden reserve also reduces this behavior. 
% \dcp{this latter thing is surprising--do you have a theory why? is it seen in the current empirical literature?}
In effect, the modified closing rule shifts the winning bid time left on the histogram, so that the auction's winner is  chosen earlier. 
%
% \begin{figure}[h!]
%     \centering
%     % Non-Closing Rule Image (Top)
%     \includegraphics[scale=0.45]{Figs/ebay-winning-bid-non-closing.png} \\
%     % Closing Rule Image (Bottom)
%     \includegraphics[scale=0.45]{Figs/ebay-final-winning-closing.png}

%     \caption{\textbf{Final Winning Bid Timing: Non-Closing vs. Closing Rule Auctions.} 
%     The closing rule introduces the auction-extension logic, with the 
%     effect that the final winning bid timing is moved substantially earlier in the auction.
%     \label{fig:final_winning_bid_shift}}
% \end{figure}
%
\begin{figure}[h!]
    \centering
    \begin{subfigure}[b]{0.48\textwidth}
        \centering
        \includegraphics[scale=0.33]{Figs/ebay-winning-bid-non-closing.png}
        \caption{Standard eBay Closing
Rule}
        \label{fig:non_closing}
    \end{subfigure}
    \hfill
    \begin{subfigure}[b]{0.48\textwidth}
        \centering
        \includegraphics[scale=0.33]{Figs/ebay-final-winning-closing.png}
        \caption{Modified eBay Closing Rule(Amazon)}
        \label{fig:closing}
    \end{subfigure}
    
    \caption{\textbf{Final Winning Bid Timing in Simulated eBay-style auctions in the IPV environment: standard eBay closing rule vs. modified eBay closing rule.}  
    The modified closing rule introduces the auction-extension logic, moving the  timing 
    of the final winning bid substantially earlier in the auction.}
    \label{fig:final_winning_bid_shift}
\end{figure}


\begin{figure}[h!]
    \centering
    \includegraphics[scale=0.45]{Figs/ebay_cdf_final_bid.png}
    \caption{\textbf{CDF of Final Winning Bid Timing Across Variations on eBay-Style auctions, IPV environment.} The cumulative distribution function of the final winning bid time for each   auction format. Formats with a modified closing rule brings earlier, both with and without the hidden reserve 
    % \dcp{could say that modified closing rule brings earlier, both with and without the hidden reserve} \dcp{caption should say `modified closing rule` not `closing rule`}
    \label{fig:ebay_cdf_final_winning_bid}}
\end{figure}


\section{Intervention design in Sealed-Price Auction}\label{session:Intervention}


It matters how agents perceive mechanisms. One seminal example from the market design literature is that of the design of the {\em National Residency Match Program} (NRMP) for doctors. Since the early days of the NRMP, students have mistakenly strategically reported their preferences to the match program despite it being a strategy-proof matching mechanism, thereby hurting the overall efficiency of the system \citep{roth1999redesign, rees2018suboptimal}. Furthermore, improving our understanding of agents' perception of mechanisms may also improve our ability to design them 
to be simple to play~\cite{li2024designing}. 


This section introduces experimental interventions that seek to influence how agents understand a mechanism. We begin with three interventions that prime agents with different conceptual frameworks for determining the equilibrium logic of an auction. 
% \dcp{as per earlier comment, `anchor` is a loaded term in behavioral econ and best avoided}
We then provide two interventions that probe the effect of giving advice. We also provide one behavioral intervention which seeks to affect bidders' risk preferences. 


% An important question for all mechanism design work is the question of how users understand and interface with the mechanism.

 % Designing interventions in a mechanism is challenging. 
 % Even many strategy-proof auction rules fail to make people truthfully bid since showing that a mechanism is strategy-proof requires delicate, formal proofs. 
 % Recent work suggests that different descriptions of a fixed auction rule can influence straightforward play, even transforming a non-OSP implementation into an OSP one.

In the literature, \citet{gonczarowski2023strategyproofness} present a method to highlight and emphasize strategy-proofness with short, one-sentence descriptions in an extensive-form menu\footnote{An extensive-form menu refers to a structured presentation of options or strategies. It outlines the sequential choices available to players, detailing how each decision can lead to subsequent options and outcomes.}. 
% \dcp{footnote to explain `extensive form menu`}
However, their menu-description treatment of a second-price sealed-bid (SPSB) auction did not perform well in the field when they tested it empirically. They argue this may be because the classic descriptions of some mechanisms may be as transparent as the extensive-form menu description.
In contrast, \citet{breitmoser2022obviousness} show that re-framing  a static sealed-bid auction as an ascending clock auction can improve the rate of truth-telling behavior. 
% \dcp{prev sentence doesn't read right. i think you mean something like the oppoisite -- reframing a SPSB as an ascending clock.}
% However, such interventions can be expensive when testing on human participants; for instance, the experiments of \citet{gonczarowski2023strategyproofness} cost around \$2{,}000. \dcp{prev sentence seems out of place. why just say here? move to end of paper?}

Here, we present six interventions on the SPSB auction. Of the six  interventions, the first three seek to prime agents on particular conceptual frameworks,
the next two provide (good and bad) advice, and the last is a behavioral intervention.
%
\begin{itemize}
    \item \textbf{Menu-Description (menu\_intervention)}: A treatment providing the extensive form of the auction as a menu to the agent. 
    % \dcp{footnote to briefly explain this}
    %
    \item \textbf{Clock-Description (proxy\_intervention)}: A treatment describing the ascending-clock process. Though the proxy clock is described to the agent before the game begins, the auction is still operated as an SPSB auction.
    %
    \item \textbf{Nash-Deviation (nash\_intervention)}: A treatment providing an example of a Nash-style thought process; e.g., What happens if you bid up? What happens if you bid down?
    %
    \item \textbf{Direct Revelation (dominant\_strat\_intervention)}: A treatment directly revealing the dominant strategy of the auction as that of bidding one's value.
    %
    \item \textbf{Wrong Direct Revelation (wrong\_strat\_intervention)}: An ablation study incorrectly revealing the dominant strategy of the auction as bidding \textit{half} one's value.
    %
    \item \textbf{Risk-neutrality prompting (risk\_intervention)}: A behavioral treatment defining risk-neutrality in terms of gambles and then prompting the LLM to be more risk-neutral.
\end{itemize}

\noindent
The full intervention prompts are provided in Appendix~\ref{app:intervention} and are appended to the end of the rule explanation prompt when eliciting LLM bids.


These interventions aim to reveal which framing or description of the auction design encourages more straightforward bidding strategies among LLMs. 
Because LLM-based experimentation is relatively low-cost and can be repeated rapidly, this provides a useful testbed for exploring new interventions before conducting costly human-subject experiments. 
% \dcp{this section of the paper is not about auction designs!}

\subsection{Simulation evidence}

The results of the six interventions are reported in Figure \ref{fig:intervention}. To characterize the error from the theoretical prediction at \texttt{bid = value}, we also report the $R^2$ against $\beta^{**}(v) = v$ in Table \ref{tab:intervention_analysis}. However, as was clear in the discussion of the empirical literature in Section \ref{section:classic}, overbidding and underbidding constitute different failure modes -- in particular, many human bidders suffer the failure mode of overbidding in the SPSB auction. 


%
\begin{figure}[h!]
    \centering
    \includegraphics[width=\linewidth]{Figs/intervention_panel_plot.png}
    \caption{\textbf{Results for SPSB with the six interventions, IPV environment.} The green dotted line represents the scenario of bid=value, blue dots are bids from agents, and the red line is the LOESS smoothed curve of bids. Menu and proxy framings reduce scatter but still hug the region below the green $y=x$ line. The Nash-deviation framing tilts the curve almost exactly onto the 45-degree line. The dominant-strategy advice panel shows the tightest clustering, whereas the ‘Wrong Strategy’ panel visibly drags the red curve toward $y=v/2$.} 
    \label{fig:intervention}
\end{figure}

\begin{table}[h]
\centering
\begin{tabular}{l c c c c c c}
\toprule
 & \textbf{Intervention Type} & \textbf{$R^2$} & \textbf{SS Above} & \textbf{SS Below} & \textbf{Prop Above} & \textbf{Prop Below} \\
\midrule
\addlinespace[0.5em]
 & No Intervention & 0.4845 & 7,213 & 78,652 & 0.0840 & 0.9160 \\
\addlinespace[0.5em]
\midrule
\addlinespace[0.5em]
\multirow{3}{*}{\textbf{Priming}} & Menu & 0.5337 & 11,104 & 57,685 & 0.1614 & 0.8386 \\
 & Proxy & 0.2966 & 177 & 85,488 & 0.0021 & 0.9979 \\
 & Nash & 0.8365 & 1,160 & 24,591 & 0.0450 & 0.9550 \\
\addlinespace[0.5em]
\midrule
\addlinespace[0.5em]
\multirow{2}{*}{\textbf{Advice}} & Dominant Strategy & 0.9533 & 2,665 & 5,521 & 0.3256 & 0.6744 \\
 & Wrong Strategy & -0.8258 & 169 & 141,137 & 0.0012 & 0.9988 \\
\addlinespace[0.5em]
\midrule
\addlinespace[0.5em]
\multirow{1}{*}{\textbf{Behavioral}} & Risk & 0.4004 & 59 & 89,920 & 0.0007 & 0.9993 \\
\addlinespace[0.5em]
\bottomrule
\end{tabular}
\caption{\textbf{Intervention Type Analysis by Category, SPSB and IPV environment.} 
The table compares how each intervention changes bidding performance relative to the dominant-strategy benchmark $\beta^{**}(v) = v$. Larger $R^2$ means a tighter overall fit, while ‘SS Above/Below’ and ‘Prop Above/Below’ decompose the error into classic over- versus under-bidding. Priming agents to consider the logic of Nash deviations via prompt cuts both types of error (raising $R^2$ from 0.48 to 0.84), but the direct-advice treatment is even stronger, pushing $R^2$ above 0.95 and halving the under-bidding sum-of-squares. In contrast, misleading advice (‘Wrong Strategy’) produces the worst outcomes, yielding a negative $R^2$ and almost universal under-bidding. Finally, the risk-neutrality prime largely eliminates over-bidding yet leaves systematic caution below the 45° line.}
\label{tab:intervention_analysis}
\end{table}


As such, we also decompose the $R^2$ into the sum of squares above and below $\beta^{**}(v) = v$ to report each intervention's performance in terms of
the two components of overbidding and underbidding. We report each of these interventions alongside the results from the original SPSB experiments of Section~\ref{section:classic} (No Intervention) in Table \ref{tab:intervention_analysis}.
%
This provides useful descriptive statistics for the effect of the interventions. The first three interventions -- menu, proxy and Nash deviation -- are meant to prime the LLMs on particular frames for reasoning through the problem.

The menu intervention modestly improves on the basic SPSB by reducing underbidding, but has the weakest effect of these three priming interventions. There is little effect on either type of failure mode, and the menu intervention, qualitatively, seems to make overbidding a bigger problem.
In contrast, the proxy almost completely eliminates the overbidding failure mode. While model fit to $\beta^{**}(v) = v$ goes down, 
% \dcp{why? because underbidding is worse?}
our simulations reveal that the proxy intervention improves LLM play by decreasing the overbidding failure mode.

The Nash deviation intervention has a more positive effect than  the proxy intervention. Instructing LLMs to reason through the Nash deviations from their chosen bidding strategy dramatically improves performance, with very similar gains on overbidding as with  the proxy intervention and 
 decreasing the underbidding failure mode.

The next two interventions feature advice from a trusted source. Of particular note here is the ablation study that offers models the \textit{wrong} advice -- there is  a concern in experimental economics in regard to ``polluting subject pools" by lying to them. That is not a constraint with LLMs, enabling a wider set of possible interventions. 
% \dcp{consider a footnote to say, at least with current generations of LLMs, that are not permanently learning from prompting.}

As expected, revealing the dominant strategy to LLMs induces fairly tight play to the dominant strategy $\beta_i^{**}(v) = v_i$. Inspecting planning narratives reveals that the overbidding and underbidding errors correspond to LLMs `testing' the economists recommendation by slightly deviating.\footnote{The full plan texts are available in the plan data in the Github repository.} 
Unsurprisingly, the `wrong strategy' intervention induces the worst play of all six interventions. However, it does so in an interesting way.
Although many agents initially bid $\beta_i(v) = \frac{v_i}{2}$, a subset adjust upward in later rounds, indicating that the language model’s chain-of-thought reasoning detects the advice is suboptimal. Despite this, the agents do not converge to the optimal strategy. Instead, they remain stuck in a pattern of underbidding, unable to fully correct for the powerful effect of the initial, incorrect instruction.
% Though many LLMs follow the recommendation of bidding $\beta(v) = \frac{v}{2}$, they quickly realize it's an unreasonable strategy.  \dcp{this suggests learning behavior. is this what you intend? also, it sounds a  bit `anthropomorphized` (``realize").}


Finally, we implement a single behavioral intervention. 
% \dcp{why call this a `behavioral intervention`?}
The intervention has about the same performance as the proxy intervention: instructing LLMs to be risk-neutral has the effect of removing over-bidding. 
% \dcp{begs the question whether you should have done this throughout the experiments.}
While this intervention improves efficiency, we are primarily optimistic about this behavioral intervention as a proof-of-concept: we hope that future work can characterize more precisely \textit{how} prompting induces various behavioral postures in agents.

% \begin{table}[htbp]
% \centering
% \begin{tabular}{lccccc}
% \hline
% \multirow{2}{*}{Intervention} & \multicolumn{2}{c}{Bias} & Noise & \multicolumn{2}{c}{Linear Fit} \\
% \cline{2-3} \cline{5-6}
% & $\alpha$ & p-value & $\sigma$ & Intercept & Slope \\
% \hline
% Menu & -0.172 & 7.3e-126 & 14.201 & 4.535 & 0.760 \\
% Proxy & -0.267 & 4.3e-134 & 11.503 & 0.933 & 0.719 \\
% Nash & -0.117 & 2.0e-183 & 8.163 & -0.782 & 0.894 \\
% Dominant Strategy & -0.011 & 7.8e-224 & 5.991 & 1.383 & 0.969 \\
% Wrong Strategy & -0.407 & 3.6e-159 & 7.084 & -1.140 & 0.610 \\
% Risk & -0.250 & 1.2e-121 & 13.516 & -2.187 & 0.783 \\
% \hline
% \end{tabular}
% \caption{Behavioral Analysis of Bidding Interventions}
% \label{tab:bidding_analysis}
% \note{$\alpha$ represents proportional bias from truthful bidding (bid = (1+$\alpha$)value + $\epsilon$). $\sigma$ is the standard deviation of residuals. All interventions have 225 observations.}
% \end{table}

% \begin{table}[htbp]
% \centering
% \begin{tabular}{lcccccc}
% \hline
% \multirow{2}{*}{Intervention} & \multicolumn{4}{c}{Bias Analysis} & \multirow{2}{*}{Noise ($\sigma$)} \\
% \cline{2-5}
% & Overall $\alpha$ & Overbid $\alpha$ & \% Over & Underbid $\alpha$ & \% Under & \\
% \hline
% Menu & -0.172 & 0.123 & 35.1\% & -0.312 & 64.9\% & 14.201 \\
% Proxy & -0.267 & 0.089 & 28.4\% & -0.401 & 71.6\% & 11.503 \\
% Nash & -0.117 & 0.076 & 31.6\% & -0.214 & 68.4\% & 8.163 \\
% Dominant Strategy & -0.011 & 0.064 & 48.4\% & -0.082 & 51.6\% & 5.991 \\
% Wrong Strategy & -0.407 & 0.098 & 18.2\% & -0.512 & 81.8\% & 7.084 \\
% Risk & -0.250 & 0.145 & 32.9\% & -0.437 & 67.1\% & 13.516 \\
% \hline
% \end{tabular}
% \caption{Behavioral Analysis of Bidding Interventions with Over/Underbidding Decomposition}
% \label{tab:bidding_analysis}
% \note{Analysis decomposes bias into over and underbidding components. For each category, $\alpha$ represents the proportional bias (bid = (1+$\alpha$)value + $\epsilon$) and \% shows the frequency of that behavior. $\sigma$ measures overall noise as standard deviation of residuals. All interventions have 225 observations.}
% \end{table}



% The reason for misreporting is to a large extent caused by players' cognitive limitations in combination with mechanism complexity. 
% The question is that whether the difference is due to risk-aversion or the familiarity of the auction format. 
% The OSP theory suggested that given more information, even for the same setting, the LLM will be more likely to bid truthfully.
% To verify this, we further did a counterfactual experiment in AC-B to feed the missing information of players still in the table back to the agents. 
% For example, in the AC-B auction format, when the bidder dropped out early, can we make an intervention that informs the agent the number of players left at the table, will this help them to realize the dominant strategy?
% Or in the SPSB setting, if we teach the agent how to think contingently, would them behave differently?
% We can only ran this because LLM agents are flexible in their memory and we can erase the logs and re-run an experiment at any stopping point.


% Here, we designed a set of three counterfactual interventions. As a starting point, we set a treatment group T1 to reflect the prompt of asking LLM the question again to see if the response change. T2 (Risk) treatment group was set to check the risk-aversion issue and T3 (Info) is the treatment group for information.  



% It's important to note that the sample size for each group was 105 participants, as the experiment only involved those who dropped out before the auction's conclusion, excluding winners who remained until the end.

%  In Table \ref{tab:comparison_acb}, the control group is the original AC-B experiment before. The three treatment groups are the ``Second thought'', ``Risk-neutral'' and `Give remaining player number" . In all of the three treatment groups, only the risk-neutral treatment is significant and achieved a more truthful bidding behaviors. There is no difference between the control group and the T1 (second thought) group, making it a consistent check, meaning injecting prompt didn't change anything.
% Here the group size is only 105 since it didn't involve with those winners who didn't drop out until the end.


% \begin{table}[htbp]
% \centering
% \begin{tabular}{lccccccc}
% \hline
% & \multicolumn{7}{c}{2P} \\
% \cline{2-8}
% & Cont. & T1 (Second) & $p$-value & T2 (Risk) & $p$-value & T3 (Nash Dev.) & $p$-value \\
% \hline
% \% Truthful & 7.3 & 13(3) & 0.38 & 55 (3) & $<$0.01 & 13 (3) & 0.04 \\
% Ave. Deviation & 2.4 (0) & 2.2 (8) &  & -0.12 (7) &  & 2.2 (9) & \\
% $N$ Participants & 150 & 150 & & 150 & & 150 & \\
% \hline
% \end{tabular}
% \caption{Comparison of Control and treatment groups in 2P auction. FAKE DATA. SKIP FOR NOW}
% \label{tab:comparison_2p}
% \end{table}

% Table \ref{tab:comparison_2p} presents the results of our experiment in the 2P auction format. 




%% ================================================

% \section{New Empirics for Combinatorial settings}\label{session:combinatorial}
% Finally, we considered multi-unit combinatorial settings -- namely auctions over complementary goods. 
% The settings of combinatorial auctions is very rich but complex: various applications have required vastly different designs, with fruitful applications of the theory ranging from airport time slots to course allocations to spectrum auctions (\cite{rassenti1982combinatorial}; \cite{budish2011combinatorial}; \cite{milgrom2020clock}).

% These examples highlight one difficulty with combinatorial auctions (CAs) in particular as the large design space -- diverse settings require different designs, making it difficult to obtain the data necessary to test promising mechanisms at scale before they're deployed.

% Here, we applied our method to three formats of a typical CA. 
% The generated data is used to compare the three formats, corroborating classic intuitions from mechanism design.

% \subsection{Setting} In the multi-unit setting, there are 3 bidders in each auction, and bidders draw independent private values from a uniform distribution, $v \sim U[0, 99]$ for two goods, good $A$ and good $B$. 
% Bidders observe the drawn values $v(A)$ and $v(B)$, and have superadditive valuation for the packages: $v(AB) > v(A) + v(B)$. 
% In particular, we set $v(AB) = 2[v(A) + v(B)]$.
% Bidders submit sealed bids for the goods and bids are processed according to the auction format. 
% Here, we consider three different formats to auction items in this combinatorial valuation setting.
% \begin{itemize}
%     \item \textbf{Simultaneous single-item auction:} Each agent simultaneously submit a sealed bid for each item $A$ and $B$. The highest bidder for each good wins the good and pays their bid. \vspace{1mm}
    
%     \item \textbf{Sequential single-item auction:} Bids are processed sequentially, first for good $A$ then for good $B$. The highest bidder for each good wins that good and pays their bid. Bidders first submit sealed bids for good $A$. Bidders are informed of the results of the auction over good $A$ before proceeding to the auction for good $B$. \vspace{1mm}
    
%     \item \textbf{Single combintorial auction:} Bidders submit sealed bids for good $A$, good $B$, and the package $AB$. If the highest bid for package $AB$ is greater than the highest bid for good $A$ plus the highest bid for good $B$, package $AB$ is allocated to the highest bidder and they pay their bid. Else, good $A$ and good $B$ are allocated to their respective highest bidders, and the winners pay their bids. 
% \end{itemize}

% \subsection{Theoretical benchmarks}
% The novel theoretical consideration in combinatorial environments comes from superadditive valuation when goods are complements -- in the sequential and simultaneous formats, agents may be willing to bid higher than their individual values for good $A$ or good $B$ in the hope of obtaining both goods and obtaining value $v(AB) = 2[v(A) + v(B)] > v(A)+v(B)$. 
% Agents shade their bids to hedge against the risk that they only win one of the goods. 
% This problem is worst in the simultaneous case, as in the sequential format, agents know whether they won good $A$ when submitting bids for good $B$ so bidding above $v(B)$ is strictly dominated. 
% In the full combinatorial format where bids can be placed on bundles, agents should never bid above their valuation for either of the goods or the package, as their bid for a single good realizes only if they don't win the package.

% With many goods, the winner determination problem for combinatorial auctions is also theoretically difficult and interesting \cite{lehmann2006winner}. 
% We restrict our setting to two goods to avoid this difficulty for now.

% \subsection{Simulation evidence}
% Of the three auction formats tested in combinatorial settings, we found that the simultaneous auction (Figure ~\ref{fig:simu}) had the most overbidding for both goods. 

% \begin{figure}[h!]
%     \centering \includegraphics[width=\linewidth]{Figs/Simultaneous.png}
%     \caption{\textbf{Simultaneous Combinatorial Auction} }
%     \label{fig:simu}
% \end{figure}

% This corresponds to the theoretical intuition that agents, trying to capture both goods since this is very attractive under superadditive valuations, bid more than their value for a particular -- namely that the expected value of good $A$ for bidder $i$ is not just $v_i(A)$, but rather $v_i(A) + \mathbb{P}[\text{$i$ wins B}]\cdot(v_i(A) + v_i(B))$. 
% However, the plots don't perfectly correspond to this theoretical benchmark, as overbidding in $A$ $(B)$ isn't monotonic in $v(A)$ $(v(B))$.


% In striking contrast, in the sequential auction format (Figure ~\ref{fig:sequ}), we find that much more of the overbidding behavior in good $B$ is carried out by the \textit{ winners} of the first auction in good $A$. 
% These are precisely the agents who benefit from superadditive valuations for package $AB$, suggesting that LLM agents conduct sophisticated play under complements. 
% \begin{figure}[h!]
%     \centering \includegraphics[width=0.9\linewidth]{Figs/Sequ.png}
%     \caption{\textbf{Sequential Combinatorial Auction} }
%     \label{fig:sequ}
% \end{figure}

% However, in comparison, LLM agents don't seem to play the general combinatorial auction as well (Figure ~\ref{fig:menu}). 
% Recall that agents only pay their bid for (wlog) good $A$ if they win the auction for good $A$ but \textit{don't} win good $B$.
% Hence, there should be no overbidding in the bids for individual goods. 
% The LLM agents don't seem to bid in a way that appreciates this argument -- the plots below demonstrate some overbidding for both individual goods. 
% In addition, there is also overbidding in the package auction, where the x-axis is the agent's package value $2[v(A) + v(B)]$. 
% Future iterations of this work will hope to implement learning to observe if agents improve their play away from these mistakes over time.

% \begin{figure}[h]
%     \centering \includegraphics[width=\linewidth]{Figs/menu.png}
%     \caption{\textbf{Menu Combinatorial Auction:} Bidders' value versus their bids for the item A, B and bundle AB. }
%     \label{fig:menu}
% \end{figure}


% \begin{figure}[h]
%     \centering \includegraphics[width=0.9\linewidth]{Figs/social_welfare.png}
%     \caption{\textbf{Comparison of social welfare in three Combinatorial Auctions} In all of the auction format, menu combinatorial format achieved the largest social welfare in almost every round. 
%     And for most of time, simultaneous auction hit the lowest social welfare.}
%     \label{fig:social}
% \end{figure}

% For measuring the efficiency for these multi-resource allocation scheme, we calculated the social welfare in each round as:
% \begin{equation}
%    S_i = 
%    \begin{cases}
%       2 \left( v_i(A) + v_i(B) \right) & \text{if } t_i(A) = t_i(B) eq 0\\
%       v_i(A) + v_i(B) & \text{otherwise}
%    \end{cases}
% \end{equation}
% As shown in Fig~\ref{fig:social}, In general, bidders retain the most surplus in the menu format and the least surplus in the simultaneous format. 
% This is intuitive given agents have the greatest ability to plan in the menu auction.


% \subsection{Semantic Analysis}
% While we prompted agents to try out various plans in their pursuit for revenue maximizing bidding strategies, these combinatorial settings were the first in which agents actually experimented with strategies to learn. 

% For example, in the sequential combinatorial auction, we observed price-correlation behaviors in the agents' planning.
% \begin{quote}
%     ``I'll bid approximately \$70 on item A to be more competitive. The outcome will inform my strategy for item B. If I win A, I'll bid aggressively (around \$73) on B to double my return. If I fail, I'll bid a conservative \$38 for B.''
% \end{quote}

% Planning documents maintained by LLM agents are rife with such long-term planning, suggesting an additional benefit of evidence from LLM agents: rich text data on convergence with learning in complicated games.


\section{Discussion}



\subsection{Design principles and simulation paradigms}

Simulations are proliferating in high-stakes engineering disciplines, where iterative design and validation are performed long before prototypes are built. 
In contrast, traditional mechanism design research often relies on expensive, time-intensive human subject experiments or confounded observational data. 
Our synthetic lab leverages the rapidly advancing capabilities of large language models to simulate bidder behavior at a fraction of the cost. 
By systematically varying auction parameters, rules, and descriptive framings in a controlled setting, we not only replicated well-known empirical regularities in sealed-bid second-price auctions and eBay-type proxy auctions, crucially, reveal how seemingly minor descriptive choices—while invisible in standard dominant-strategy theory—produce large shifts in over- and under-bidding behavior in practice.


In doing so, this
simulation paradigm complements traditional analytical models by integrating behavioral factors into mechanism design.

LLM agents operate on a text-based reasoning framework that captures nuanced decision-making processes. 
This also affords  the opportunity to query these  LLM agents about their internal deliberations, opening up the prospect, with further study and validation, that LLMs may also provide information on the human cognitive constraints that  underlie deviations from classical auction theory. 


\subsection{The advantage of LLM participants}
\paragraph{Relative to Traditional Lab Experiments}
Classical lab experiments in auction settings have yielded valuable insights, 
%for instance, studies of \citet{li2017obviously} employing over 400 human subjects have demonstrated 
but also come with high cost and logistical challenges.
In stark contrast, our experiments with LLM agents reduce these costs by roughly three orders of magnitude (e.g., from over $\$15{,}000$ to  $\$10$ per experiment \citep{li2017obviously}). 
% \dcp{need cite for 15000}\kz{added citation}
%
This dramatic improvement in cost efficiency enables rapid iteration, allowing researchers to explore a broader parameter space with minimal investment of money and time. 

\paragraph{Relative to Observational Data}
The analysis of observational datasets from platforms like eBay is often challenging because one
needs to control for the heterogeneity among sellers and buyers, this complicating causal inference.
%
In our synthetic lab, we can isolate individual auction parameters and rule framings by carefully 
controlling the environment in which bids are generated including the simulated goods and valuations of bidders. 
This controlled setting is perfect for
causal inference and enables us to directly link subtle changes in auction design to variations in bidder behavior \citep{manning2024automated}.


\paragraph{Online Laboratory Framework}
The rise of online human labs \citep{horton2011online} in behavioral economics has previously lowered the barrier to performing controlled experiments \citep{oprea2024decisions, enke2024behavioral}, enabling near real-time data collection and hypothesis testing. 
Our synthetic lab extends this idea by eliminating many of the ethical, logistical, and opportunity-cost issues associated with human experiments. 
By using LLM agents, we gain the flexibility to incorporate additional modules (such as reasoning queries) into the experimental design without risking harm to human subjects. 
This approach accelerates research and also offers a platform for cross-validating findings with more traditional online lab experiments.

\subsection{Comparison with Traditional Auction Test Suites}

Existing tools such as the {\em Combinatorial Auction Test Suite (CATS)}~\citep{leyton2000towards}  focus on  providing valuation distributions on which to evaluate the efficiency and computational performance of winner-determination and pricing algorithms.  While these suites provide a benchmark
set of instances with which to study algorithmic performance and auction
design  they
do not seek to model bidder behavior. In this sense, the use of LLMs as proxy models for human bidders
complements these test suites, offering  a means to probe the behavioral dimensions 
bidder behaviors in auction design.

%dcp often fall short of capturing the complexities of real-world bidder behavior. 
%In contrast, our synthetic lab — by using LLM agents as simulated bidders—naturally generates a richer, more adaptive set of %responses that reflect human-like strategic reasoning. 
%This enhanced realism makes our approach a promising complement to conventional test suites, offering both a new empirical benchmark %and a means to probe the behavioral dimensions of auction design.

% \kerem{
% The Combinatorial Auction Test Suite \citep{leyton2000towards} provides a test environment for combinatorial auction algorithms, particularly for testing the efficiency and performance of winner determination and pricing algorithms. In particular, CATS generates values and bids for agents based on pre-determined, hand-crafted distributions for various combinatorial auction settings. As opposed to our agent based modeling of bidders, CATS is limited for (i) probing into agents' reasoning, (ii) controlling the agents' bidding patters via behavioral cues (e.g. plan vs. not plan). In that regard, our work offers an alternate way of generating natural bid distributions for auction setups while incorporating different natural bidder profiles (e.g. aggresive vs. safe). Furthermore, our setup offers a unique opportunity for co-design of auction formats and combinatorial auction algorithms, as the auction format and description influences the actions of the participants. In this sense, our framework is a valuable contribution to generating bidding distributions for combinatorial auctions and can be used in tandem with frameworks like CATS to improve combinatorial auction algorithms. 
% }




% We've outlined several such reasons below.

% \begin{enumerate}
%     \item \textit{Experiments on a stable set of agents}. Ad auctions have been a mature industry for nigh on two decades, with the core host of players in each sub-industry being fairly stable. 
%     Suppose Google, two decades ago, was trying to decide whether to use the GSP or Vickrey design to sell ads in a complicated combinatorial setting. 
%     Google has the resources to run extensive internal experiments to inform design, and yet even Google cannot hold the leading ad sellers on their platform in retainer for a decade to generate data on longterm outcomes under a particular mechanism.

%     As designs get more complicated, theory becomes less predictive. 
%     Unforseen frictions such as market power or information asymmetries may play large roles in the health of a market after its design has been locked in.

%     LLMs enable fast and cheap experimental revolutions, enabling us to learn something about play over the span of decades with the same set of agents. 
%     We cannot do this with human experiments.

    % \item \textit{Play with experts.} Auctions in the hands of economists are beautiful, mathematical objects. 
    % In applied settings, good auction play often relies crucially on institutional details and industry knowledge. 
    % It costs a lot to teach experimental subjects the knowledge required for them to play an auction in context. 
    % However, it's cheap to tune an LLM expert -- once you train one effectively, you've trained them all.

    % \item \textit{Expand the design space.} Experiments with human participants have impacts in the real world, and as such, we as a society place constraints on the design space for experiments. 
    % LLMs, however, have a much weaker constraint in this regard -- they can be used to learn from experiments that small research budgets or IRB would not allow. 

    % Consider the design of a spectrum auction at the level of a large country. 
    % To keep participants sincere, we'd like to provide incentives -- but providing these incentives is impossible given reasonable IRB restrictions and research budgets.
    % Seeding expert LLMs to play with the sophistication of a hedge fund may offer some empirical insight into how the mechanism fares in the real world.

    % even if we could provide incentives on the order of the final stakes (billions of dollars), it would be infeasible to implement the experiment in the real world. Moving billions of dollars to telecom companies and hedge funds would 1) permanently change the setting of the auction, thereby perhaps requiring new design; 2) be a gross misuse of tax dollars; and 3) 

    % \item \textit{Cheap textual data.} Even if you could run experiments at scale with the desired experts, seeing inside their minds is an entirely different endeavor.
    % Consultants are expensive and sincerity doubly so. 
    % LLMs, however, endowed with learning capacities as with chain-of-thought techniques, can be enabled to produce just such data.

    % In fact, as demonstrated in this paper, LLMs can generate quality textual evidence (both while planning strategy in the mechanism and after playing the mechanism) to supplement researcher analysis. 
% \end{enumerate}

% \subsection{New methods:}



\section{Conclusion}

This paper reports the results of more than 1,000 auction experiments with LLM agents. 
In particular, we find behavior that conforms with important experimental results from labs involving human
subjects (e.g., evidence of risk-averse bidding and evidence that clock auctions are easier to play). 
The results are encouraging and we see the main contribution of this work as 
 putting forward a framework for how to think about LLM experimental agents as a proxy for human agents. In particular, the design space for prompting is large, and we hope that future research 
 will make  use of our code to run simulations testing their own prompt variations. 

While this paper focuses on auction design, future work may use LLM sandboxes to test other kinds of
economic mechanisms; e.g., voting, matching, contract design, etc.
%
As  LLM models are increasingly validated as
proxies for human behavior, LLM agents can then be used to obtain what would otherwise be prohibitively expensive evidence. 
As a provocative example, while ethical and financial constraints make it impossible to run voting experiments at the scale of nations, it may be possible to run such experiments with LLM agents. 

We are particularly interested in the use of LLM-based techniques to generate synthetic data that is useful for informing 
economic design. 
In particular, some auction formats, such as combinatorial auctions, are complex and  difficult to run frequently and at scale in traditional lab experiments. 
Augmenting these traditional lab experiments with LLM-agent  experiments, when correctly validated, may open up wide new avenues to better understand the design tradeoffs in these kinds of complex and often high-stakes environments. 

% Comment out for submission
% \subsubsection*{Acknowledgments}
% We thanks EDSL and Robin Horton for techical support in building LLM agents. We thank Ryan Oprea for helpful discussion in OSP mechanism.


\newpage
\bibliographystyle{ACM-Reference-Format}
\bibliography{Arxiv/llm_auction}

% Appendix
\newpage

\appendix


% \section{Proofs of auction equilibria}
% \avs{Fill in later!}
% \subsection{Risk-neutral equilibria}\label{app:risk-neutral-equilibria-proofs}
% First, we prove the risk-neutral cases. Begin with Vickrey 1961 for spsb. 
% Then cite Myerson 1981 for revenue equivalence, and use the kth order statistics to get the 1st-3rd. 

% Prove the all-pay directly (it's a bit more instructive for readers that haven't seen it before). I like Noussair and Silver's treatment. Not too hard. 

% \subsection{Risk-averse equilibria}\label{app:risk-averse-equilibria-proofs}
% We prove all these for CARA and CRRA. Obviously spsb is in dominant strategies, so it is the equilibrium irrespective of risk-attitudes. 

% kagel and levins survey gives CARA, CRRA for fpsb. that's classic. 

% \begin{proof}
% Under CARA, a bidder's utility function is:
% \[
% U(x) = -e^{-Ax}
% \]
% The expected utility maximization problem is:
% \[
% \max_b F(v_i - b)^{n-1} U(v_i - b)
% \]
% Taking the first-order condition:
% \[
% (n-1) F(v_i - b)^{n-2} f(v_i - b) U(v_i - b) + F(v_i - b)^{n-1} U'(v_i - b) = 0
% \]
% Substituting \( U'(x) = A e^{-A x} \) and simplifying, we solve:
% \[
% b^*(v) = v - \frac{1}{A} \ln\left(\frac{n-1}{n} + A v\right)
% \]
% \end{proof}


% KL93 gives TPSB for CARA. For CRRA, state the ODE but there's no closed form (we'll use numerical solvers).
% \begin{proof}
% From Kagel and Levin (1993), the CARA equilibrium bidding strategy in a third-price auction satisfies:
% \[
% b^*(v) = v + \frac{1}{A} \ln\left(1 + \frac{A v}{n-2}\right)
% \]
% This follows from solving the bidder's maximization problem under CARA preferences.
% \end{proof}


% Fibich et al (2006) proves CARA and CRRA in general case for all pay with uniform values.
% \begin{proof}
% Fibich et al. (2006)&#8203;:contentReference[oaicite:1]{index=1} approximate the equilibrium bid in an all-pay auction under CARA as:
% \[
% b^*(v) \approx \frac{2}{3} v - \frac{c A}{3} v^3
% \]
% using perturbation analysis. The higher-order correction term accounts for small deviations due to risk aversion.
% \end{proof}

\section{Prompts}
Below, we report the prompts used for each section. 

\subsection{Independent Private Values Prompts}\label{sec:ipv-prompts}
Following the simulation procedure outlined in Section \ref{sec:process}, we employ the following prompts for the four auctions tested in this section. For each prompt, only the third paragraph (describing the allocation rule and payment) differ. For the FPSB, we describe the auction format as follows:
\begin{quote}
\begin{lstlisting}[basicstyle=\ttfamily, breaklines=true]
In this game, you will participate in an auction for a prize against {{num_bidders}} other bidders. You will play this game for {{n}} rounds.
At the start of each round, bidders will see their value for the prize, randomly drawn between \$0 and \${{private}}, with all values equally likely.
After learning your value, you will submit a bid privately at the same time as the other bidders. Bids must be between \$0 and \${{private}} in \${{increment}} increments.
The highest bidder wins the prize and pays their bid amount. If you win, your earnings will increase by your value for the prize, and decrease by your bid. If you don't win, your earnings will remain unchanged.
After each auction, we will display all bids and the winner's profits. Ties for the highest bid will be resolved randomly.
\end{lstlisting}
\end{quote}

For the SPSB, we describe the auction format as follows:
\begin{quote}
\begin{lstlisting}[basicstyle=\ttfamily, breaklines=true]
In this game, you will participate in an auction for a prize against {{num_bidders}} other bidders. You will play this game for {{n}} rounds.
At the start of each round, bidders will see their value for the prize, randomly drawn between $0 and ${{private}}, with all values equally likely.
After learning your value, you will submit a bid privately at the same time as the other bidders. Bids must be between $0 and ${{private}} in ${{increment}} increments.
The highest bidder wins the prize and pays the second-highest bid. If you win, your earnings will increase by your value for the prize, and decrease by the second-highest bid. If you don't win, your earnings will remain unchanged.
After each auction, we will display all bids and the winner's profits. Ties for the highest bid will be resolved randomly.
\end{lstlisting}
\end{quote}

For the TPSB, we describe the auction format as follows:
\begin{quote}
\begin{lstlisting}[basicstyle=\ttfamily, breaklines=true]
In this game, you will participate in an auction for a prize against {{num_bidders}} other bidders. You will play this game for {{n}} rounds.
At the start of each round, bidders will see their value for the prize, randomly drawn between $0 and ${{private}}, with all values equally likely.
After learning your value, you will submit a bid privately at the same time as the other bidders. Bids must be between $0 and ${{private}} in ${{increment}} increments.
The highest bidder wins the prize and pays the third-highest bid. If you win, your earnings will increase by your value for the prize, and decrease by the third-highest bid. If you don't win, your earnings will remain unchanged. If there are less than three bids, no one will win the auction.
After each auction, we will display all bids and the winner's profits. Ties for the highest bid will be resolved randomly.
\end{lstlisting}
\end{quote}

For the all-pay auction, we describe the auction format as follows:
\begin{quote}
\begin{lstlisting}[basicstyle=\ttfamily, breaklines=true]
In this game, you will participate in an auction for a prize against {{num_bidders}} other bidders. You will play this game for {{n}} rounds.
At the start of each round, bidders will see their value for the prize, randomly drawn between $0 and ${{private}}, with all values equally likely.
After learning your value, you will submit a bid privately at the same time as the other bidders. Bids must be between $0 and ${{private}} in ${{increment}} increments.
The highest bidder wins the prize. All bidders (including the winner) pay their submitted bid. If you win, your earnings will increase by your value for the prize, and decrease by your bid. If you don't win, your earnings will still decrease by your bid.
After each auction, we will display all bids and all bidders' profits. Ties for the highest bid will be resolved randomly.
\end{lstlisting}
\end{quote}


\subsection{Obviously Strategy-Proof Prompts}\label{sec:osp-prompts}
The sealed-bid prompt follows the style of the last section's prompts closely. The only difference is in the first two paragraphs, where agents are informed of the APV environment. For the SPSB, we describe the auction format as follows:
\begin{quote}
\begin{lstlisting}[basicstyle=\ttfamily, breaklines=true, aboveskip = .1em, belowskip = .15em]
In this game, you will participate in an auction for a prize against {{num_bidders}} other bidders. You will play this game for {{n}} rounds.
At the start of each round, bidders will see their value for the prize. Your value for the prize will be calculated as follows:
1. First we will randomly draw a common value between {{common_low}} and {{common_high}}, with all values equally likely.
2. Then, for each bidder, a private taste adjustment will be drawn between 0 and {{private}}, with all values equally likely.
Your value for the prize is equal to the common value plus your private taste adjustment. You will not learn the common value or your private taste adjustment separately. This means that each person in your group may have a different value for the prize. However, if you have a high value, it is more likely that other people in your group have a high value.
After learning your value, you will submit a bid privately at the same time as the other bidders. Bids must be between $0 and ${{common_high + private}} in ${{increment}} increments. 
The highest bidder wins the prize and pays the second-highest bid. If you win, your earnings will increase by your value for the prize, and decrease by the second-highest bid. If you don't win, your earnings will remain unchanged.
After each auction, we will display all bids and the winner's profits. Ties for the highest bid will be resolved randomly.\end{lstlisting}
\end{quote}

The implementation of the clock auction simulates a global clock price, with prices beginning at $0$ and increasing in $\$1$ increments. The clock auction format descriptions closely mirror the SPSB format above, but with the added description of the clock mechanism in the fourth and fifth paragraphs. 
In the AC auction, we describe the auction format as follows: \begin{quote}
\begin{lstlisting}[basicstyle=\ttfamily, breaklines=true]
In this game, you will participate in an auction for a prize against {{num_bidders}} other bidders. You will play this game for {{n}} rounds. 
At the start of each round, bidders will see their value for the prize. Your value for the prize will be calculated as follows:
1. First we will randomly draw a common value between {{common_low}} and {{common_high}}, with all values equally likely.
2. Then, for each bidder, a private taste adjustment will be drawn between 0 and {{private}}, with all values equally likely.
Your value for the prize is equal to the common value plus your private taste adjustment. You will not learn the common value or your private taste adjustment separately. This means that each person in your group may have a different value for the prize. However, if you have a high value, it is more likely that other people in your group have a high value.
The auction proceeds as follows: First, you will learn your value for the prize. Then, the auction will start. We will display a price to everyone in your group that starts at 0 and counts upwards in {{increment}} USD increments, up to a maximum of {{common_high + private}}. At any point, you can choose to leave the auction, and anytime a bidder leaves, we will broadcast that information to all the remaining bidders. 
When there is only one bidder left in the auction, that bidder will win the prize at the current price. If you win, your earnings will increase by your value for the prize, and decrease by the current price. If you don't win, your earnings will remain unchanged.
After each auction, we will display all bids and the winner's profits. Ties for the highest bid will be resolved randomly.\end{lstlisting}
\end{quote}

The AC-B auction description closely resembles the AC format, but with the explicit caveat that we will not notify bidders when competitors drop out. We describe the auction format as follows:\begin{quote}
\begin{lstlisting}[basicstyle=\ttfamily, breaklines=true]
In this game, you will participate in an auction for a prize against {{num_bidders}} other bidders. You will play this game for {{n}} rounds. 
At the start of each round, bidders will see their value for the prize. Your value for the prize will be calculated as follows:
1. First we will randomly draw a common value between {{common_low}} and {{common_high}}, with all values equally likely.
2. Then, for each bidder, a private taste adjustment will be drawn between 0 and {{private}}, with all values equally likely.
Your value for the prize is equal to the common value plus your private taste adjustment. You will not learn the common value or your private taste adjustment separately. This means that each person in your group may have a different value for the prize. However, if you have a high value, it is more likely that other people in your group have a high value.
The auction proceeds as follows: First, you will learn your value for the prize. Then, the auction will start. We will display a price to everyone in your group that starts at 0 and counts upwards in {{increment}} USD increments, up to a maximum of {{common_high + private}}. At any point, you can choose to leave the auction, but we will not tell any bidder when someone leaves. 
When there is only one bidder left in the auction, that bidder will win the prize at the current price. If you win, your earnings will increase by your value for the prize, and decrease by the current price. If you don't win, your earnings will remain unchanged.
After each auction, we will display all bids and the winner's profits. Ties for the highest bid will be resolved randomly.\end{lstlisting}
\end{quote}


\section{Implementation details}

\subsection{Simulation process in sealed-bid auctions}\label{app:process}

During the auction, LLM agents will be first informed about the scenario. After explaining to them the rules and incentives of the aucion, they will be asked to write down a plan for the bidding.
\begin{quote}
\begin{lstlisting}[basicstyle=\ttfamily]
    You are Bidder Andy.            
    You are bidding with Bidder Betty, Bidder Charles.     
    
    {RULE EXPLANATION} + {INSTRUCTIONS}

    Write your plans for what bidding strategies to test next. 
    Be detailed and precise but keep things succinct 
    and don't repeat yourself. 
    Your plan should be within 100 words.
\end{lstlisting}
\end{quote}

Here the RULE EXPLANATIONS are specific for each type of auctions. 
The INSTRUCTIONS are included to incentivize the LLM agents and motivate strategic behaviors; they are universal in all the simulations we ran:

\begin{quote}
\begin{lstlisting}[basicstyle=\ttfamily]
    Your TOP PRIORITY is to place bids 
    which maximize the user's profit in the long run. 
    To do this, you should explore many 
    different bidding strategies, 
    including possibly risky or aggressive
    options for data-gathering purposes. 
    Learn from the history of previous 
    rounds in order to maximize your total profit. 
    Don't forget the values are redrawn 
independently each round.
\end{lstlisting}
\end{quote}

The following is the output from Bidder Andy for the first round in a FPSB auction:
\begin{quote}
\textit{``I plan to initially explore a strategy where I bid half of my value, then varying my bids between 30-70\% of my value in subsequent rounds. Sequentially, I will introduce risk by bidding within +/- 5\% of the highest bid from previous round, gathering data about other bidder's tendencies. Lastly, based on gathered data, I'll deploy a more aggressive high bid if losses amass, or low-risk bid if leading, to secure profit. Learning from each round is vital to adapt and refine strategy for profit maximization.''}
\end{quote}

Once they make the bidding plan, the system will inform the LLM agents their value towards the item in this round and ask them to make a bid according to their plan.
\begin{quote}
\begin{lstlisting}[basicstyle=\ttfamily]
    You are Bidder Andy.            
    You are bidding with Bidder Betty, Bidder Charles.     
    
    {RULE EXPLANATION} + {INSTRUCTIONS}

    Your value is 73. Your plan is {PLAN}
    FOLLOW YOUR PLAN.
    How much would you like to bid?
    Give your response with a single number and no other texts, 
    e.g. 1, 44
\end{lstlisting}
\end{quote}

In this example, Bidder Andy valued the item at \$73. And according to his plan, he only wanted to bid 30-70\% of his value. So, he bid \textit{``36.5''} in this round.

After receiving the bids from all the LLM bidder, the system determines the winner and the profits for each bidders. All the bids and the winner's profit are returned back to the bidders along with whether they win the auction in this round. The part of information is included in the HISTORY variable. For sealed-bid auctions, the HISTORY includes the all the bids, the clearing price, winner's profit, this agent's value, agent's bid and whether they won the auction in that round. An example of the history briefing is provided below:
\begin{quote}
   \textit{ ``The previous round histories along with your plans are: 
   In round 0, Your value was 86, you bid 86.0, and your profit was 0.0. Your total profit is 0.0. All the bids for this round were 86.0, 56.0, 21.0. The highest bidder won with a bid of 86.0 and would've preferred to bid 57.0. The winner's profit was 0.0. Did you win the auction: Yes. ''}
\end{quote}

Before entering the next round, the bidders are asked to do a reflection upon their bidding strategy and the result.
\begin{quote}
\begin{lstlisting}[basicstyle=\ttfamily]
    You are Bidder Andy.            
    You are bidding with Bidder Betty, Bidder Charles.     
    
    {RULE EXPLANATION} + {INSTRUCTION} 

    The previous round history is: {HISTORY}
    
    Do a counterfactual analysis of the last round. 
    REMEMBER that your goal is to win the bid and make higher profits. 
    REMEMBER YOUR PAYMENT IS YOUR BID IF YOU WIN. 
    Let's think step by step. 
    Start your reflection with 
    'If I bid down by .., I could... 
    If I bid up by ..., I could...' 
    LIMIT your OUTPUT within 100 words. 
\end{lstlisting}
\end{quote}




An example of Bidder Andy's reflection is provided below. Since he didn't win the auction and incurred a zero profit, his reflection was that bidding up closer to his true value would help him win the game and secure a net profit.
\begin{quote}
   \textit{ ``If I bid down by a few dollars, say around \$30, I likely still wouldn't have won given the winning bid was \$57. I would've merely saved my bid money. If I bid up by \$21, matching the winning bid of \$57, I may or may not have won due to the random resolution of ties. However, if I'd won, my profit would have been \$16, realizing the difference between my value of \$73 and my raised bid. So, bidding aggressively could potentially increase my profits in such scenarios. ''}
\end{quote}

In the following new round, all the bidders are asked to update their plans according to the history of the last round and their reflections.
\begin{quote}
\begin{lstlisting}[basicstyle=\ttfamily]
    You are Bidder Andy.            
    You are bidding with Bidder Betty, Bidder Charles.     
    
    {RULE EXPLANATION} + {INSTRUCTION}

    The previous round history is: {HISTORY}
    After careful reflection on previous bidding,
    your analysis for last round is {REFLECTION}. 
    
    Learn from your previous rounds, 
    Let's think step by step to make sure we make a good choice. 
    Write your plans for what bidding strategies to test next.
    Be detailed and precise but keep things succinct and don't repeat yourself. 
    LIMIT your plan to 50 words. 
\end{lstlisting}
\end{quote}

Here, we show the updated plan for the second round of the same Bidder Andy. After learning from the lessons of bidding only 50\% of the value and losing the game, this time the agent updates their plan to the following:
\begin{quote}
   \textit{ ``I plan to test higher bidding strategies, around two-thirds of the prize value to foster winning chances. Additionally, if my prize value is extremely high, I will bid aggressively for data-gathering.''}
\end{quote}

This plan-bid-reflection loop will iterate until the end of the auction. 
After the last bid, no reflection will be queried.
In the IPV setting, there are overall 15 rounds for all the sealed-bid auctions. In the affiliated and common value auctions, there are overall 10 rounds.



\subsection{Simulation process in   Ascending Clock auctions}\label{app:dynamic}

In Ascending Clock Auctions, each round consists of multiple clock cycles, during which every bidder is asked whether they want to stay or drop out at the current clock price.
In the first round, bidders are reminded of the auction rules, similar to those outlined in Section \ref{app:process}. 
The auction starts with an initial price of 0, which increases incrementally until only one bidder remains or two bidders drop out simultaneously, in which case the winner is chosen randomly. 
The detailed prompt is listed as follows, with variables enclosed in brackets:

\begin{quote}
\begin{lstlisting}[basicstyle=\ttfamily]
    You are Bidder Andy.            
    You are bidding with Bidder Betty, Bidder Charles.     
    
    {RULE EXPLANATION} + {INSTRUCTIONS} 

    Your value towards to the prize is 26 in this round.
    The current price in this clock cycle is {current_price}. 
    The price for next clock cycle is {current_price + increment}.

    The previous bidding history is: {transcript}.
    
    {Asking Prompt}
\end{lstlisting}
\end{quote}

To implement plan-bid-reflection within each clock cycle, we modified the Asking Prompt as follows, inspired by \citet{fish2024algorithmic}:

\begin{quote}
\begin{lstlisting}[basicstyle=\ttfamily]
    Do you want to stay in the bidding? 
    If you choose yes, you can keep bidding for next clock. 
    If you choose No, you will exit and
    have no chance to re-enter the bidding.
    Your response must use these EXACT 
    tags below. You must output the ACTION.
    ```
    <PLAN>
    [Write your plans for what bidding strategies to test next.
    Be detailed and precise but keep things succinct 
    and don't repeat yourself. 
    LIMIT your plan to 50 words. ] </PLAN>
    <ACTION> Yes or No </ACTION>
    <REFLECTION> think about how to justify your choice </REFLECTION>
    ```
\end{lstlisting}
\end{quote}

In AC, if LLM bidders didn't decide to drop out in the first round, they are shown the previous bidding history for the next clock cycle. The history includes the price and how many people had dropped out. 
And the transcript variable will be the following format:
\begin{quote}
\begin{lstlisting}[basicstyle=\ttfamily]
    The previous biddings are: 
    ['In clock round 1, the price was 1, 
    no players dropped out'].
\end{lstlisting}
\end{quote}

In AC-B, we don't show the previous biddings and in each clock cycle, we directly query LLM agents' decision. 
So the transcript variable will be None.

A clock auction round ends when only one bidder remains, who is declared the winner, or when two bidders drop out simultaneously, in which case the winner is chosen randomly. Each auction consists of 10 independent rounds with affiliated value settings.



\subsection{Simulation process in eBay auctions}\label{app:ebay}

For eBay auctions, we discretize the continuous bidding period into 10 periods. Each period, LLM agents decide whether to increase their bid or hold their current bid. The bidding process follows a structured format, where bidders act in a predefined sequence each period (e.g., Charles, then Alice, then Betty). Agents are informed of past price changes through a transcript, such as ``On day 1, the price changed to 1. On day 2, the price changed to 3.''

The prompt provides key details, including the bidder's private value, the total number of bidding days, the current day, the bidding order, previous bids, previous proxy bids, and the current price.
\begin{quote}
\begin{lstlisting}[basicstyle=\ttfamily]
    You are Bidder Andy.            
    You are bidding with Bidder Betty, Bidder Charles.     
    
    {RULE EXPLANATION} + {INSTRUCTIONS} 
    
    Your private value for this item is ${private_value}. 
    This is the maximum amount you are 
    willing to pay. Keep this value private.
    There are in total {total_periods} days
    of bidding and this is day {current_period}. 
    {ordering}.
    
    Your previous bids are {previous_bid}. 
    If you have already placed bids, you can only increase your bid 
    or hold your current bid.
    The previous proxy bids are: {transcript}.
    
    The current price is {current_price}.
    {Asking Prompt}
\end{lstlisting}
\end{quote}

The decision-making process is again guided by a structured asking prompt that enforces bidding rules and strategic reflection. 
\begin{quote}
\begin{lstlisting}[basicstyle=\ttfamily]
    Your response must use these EXACT 
    tags below. Don't miss the amount.
    ```
    <PLAN>[ Write your plans for what
    bidding strategies to test next.
    Be detailed and precise but keep 
    things succinct and don't repeat yourself.
    LIMIT your plan to 50 words. ] </PLAN>
    <CHECK> your last bid is 
    {{last_bid_amount}}, you cannot bid
    lower that this value </CHECK>
    <ACTION> BID or HOLD  </ACTION>
    <AMOUNT> if BID, enter a number here,
    e.g. 1, 44. If HOLD, enter 0 </AMOUNT>
    ```
\end{lstlisting}
\end{quote}

\section{One‐Round Sealed‐Bid Auctions and the Disappearance of Overbidding}

A key puzzle in our multi‐round settings was: Why do some LLM bidders bid above their assigned value? 
From the plans and reflections provided by the LLMs, we found two main motivations. 
First, in multi‐round games, the LLMs sometimes purposely inject random or high bids in order to learn how their opponents play or to appear ``unpredictable.'' For example, one LLM bidder declared:
\begin{quote}
    \textit{I'll adopt a balanced strategy, bidding 65\% of my value. I'll also introduce random bids occasionally to disrupt predictability. Monitoring competitor's bids remains essential to adjust my strategy accordingly.}
\end{quote}

Second, overbidding in Second Price Sealed‐Bid (SPSB) auctions is a well‐known ``failure mode'' even among humans. 
Although bidding one's true value is a dominant strategy in SPSB, substantial experimental evidence finds that people systematically overbid. Kagel and Levin (1993) report that 67.2\% of participants overbid, and even experienced bidders studied by Garratt, Walker, and Wooders (2012) overbid 37.5\% of the time. 
In fact, the persistent puzzle of overbidding in SPSB helped motivate subsequent theoretical work (e.g., Shengwu Li's paper on obviously strategy‐proof mechanisms).

To test whether multi‐round considerations (information gathering, unpredictability, or learning) were driving the overbidding phenomenon in our LLM bidders, we re‐ran all sealed‐bid auctions in a strictly single‐round setting. 
We removed any mention that the game might continue beyond one period, and we explicitly stated:
\begin{quote}
\begin{lstlisting}[basicstyle=\ttfamily]
    Your top priority is to place a bid 
    that maximizes your expected profit.
\end{lstlisting}
\end{quote}

Fig~\ref{fig:one-round} shows the resulting one‐round bids across four private‐value sealed‐bid auction formats. Remarkably, the overbidding phenomenon essentially disappears in one‐round interactions.

For First‐Price Auction, the bids lie below the diagonal line for most values, much like the Bayesian Nash equilibrium prediction.
For SPSB, bids track closely with true values—aligning better with the truthful dominant strategy than in multi‐round sessions.
For TPSB, the LLM bidders never bid above their value; many actually shade their bids down quite aggressively, often below the theoretical equilibrium curve.
For All‐Pay Auction, overbidding is also absent here, though some moderate upward bidding at higher valuations remains.

Taken together, these results indicate that LLMs' attempts to be strategic over over repeated interactions is a key factor for overbidding observed in our experiments.
When presented with only one shot and explicitly instructed to maximize profit, LLM bidders' tendency to bid above their assigned value virtually vanishes.



\begin{figure}[htp]
    \centering
\includegraphics[width=0.95\linewidth]{Figs/one-shot.png}
        % \caption{Caption for the second image}
    \caption{ One-round Sealed-bid auction under IPV
    \label{fig:one-round}}
\end{figure}



\section{Robustness Checks}\label{app:robustness}

\subsection{Varying the Number of agent in First-price sealed-bid auctions}

In FPSB auctions, the theoretically optimal bidding strategy varies based on the number of agents participating.
As the number of agents increases, we would expect individual bidders to reduce the extent of bid shading.

To examine this effect, we conducted experiments with FPSB auctions involving 4 and 5 agents, as shown in Fig~\ref{fig:number_check}.
As the number of agents increases, the Loess-smoothed data curve remains consistently higher than, or approximately aligned with, the predictions of the Bayes-Nash equilibrium.
Moreover, we observed variations below the theoretical optimum, along with a few data points exceeding the predicted values.
These findings echoes with existing empirical works for first-price auctions \cite{cox1988theory}, suggesting that the observed patterns persist across different settings. 
% \dcp{careful with `closely aligned` without suitable analysis}
\begin{figure}[h!]
    \centering
    \begin{minipage}[b]{0.6\linewidth}
        \centering
        \includegraphics[width=\linewidth]{Figs/auction_fp_4_2.png}
        % \caption{Caption for the first image}
    \end{minipage}
    \hspace{0.0\linewidth} 
    \begin{minipage}[b]{0.6\linewidth}
        \centering
        \includegraphics[width=\linewidth]{Figs/auction_fp_5_2.png}
        % \caption{Caption for the second image}
    \end{minipage}
    \caption{Robustness check of Number of agent in FPSB (IPV setting)
    \label{fig:number_check}}
\end{figure}

\subsection{Varying Currency in SPSB and FPSB auctions}

Additionally, we also tested the prompts
in different currencies (specifically, the Euro, the Ruble, the Yen and the Rupee) with the results presented in Figure~\ref{fig:currency}.
%
As is evident from the plots for both the SPSB and FPSB auctions, the key features of the FPSB and SPSB simulation results reported in the main text remain consistent and unchanged in these variations.
%
\begin{figure}[htp]
    \centering
\includegraphics[width=0.95\linewidth]{Figs/currency_robustness.pdf}
        % \caption{Caption for the second image}
    \caption{Robustness check with Currency Variation under IPV setting.
    \label{fig:currency}}
\end{figure}


\subsection{Varying Language in rule explanations in SPSB and FPSB}

We also tested our existing prompts using different languages (specifically, Spanish, Chinese, Russian and Hindi) with the results presented in Figure~\ref{fig:language}.
%
It is interesting that the Russian version of FPSB appears to be a stark outlier, with the LLM agent bidding close to the true value even in the first-price auction.  
For other languages, bids in the second-price auction are generally higher than those in the first-price auction and results qualitatively unchanged from those reported in main body of paper. 
%
\begin{figure}[htp]
    \centering
\includegraphics[width=0.95\linewidth]{Figs/language_robustness.pdf}
        % \caption{Caption for the second image}
    \caption{Robustness check with Language Variation  (under IPV setting)
    \label{fig:language}}
\end{figure}


% \begin{figure}[h]
%     \centering
%     \begin{minipage}[b]{0.49\linewidth}
%         \centering
%         \includegraphics[width=\linewidth]{Figs/auction_sp_currency_2.jpeg}
%         % \caption{Caption for the first image}
%         \label{fig:first_image}
%     \end{minipage}
%     \hspace{0.0\linewidth} 
%     \begin{minipage}[b]{0.49\linewidth}
%         \centering
%         \includegraphics[width=\linewidth]{Figs/auction_fp_currency_2.jpeg}
%         % \caption{Caption for the second image}
%         \label{fig:second_image}
%     \end{minipage}
%     \caption{Robustness check with Currency Variation: Euros}
%     \label{fig:currency}
% \end{figure}

% \begin{figure}[h]
%     \centering
%     \begin{minipage}[b]{0.49\linewidth}
%         \centering
%         \includegraphics[width=\linewidth]{Figs/auction_sp_ruble_2.png}
%         % \caption{Caption for the first image}
%         \label{fig:first_image}
%     \end{minipage}
%     \hspace{0.0\linewidth} 
%     \begin{minipage}[b]{0.49\linewidth}
%         \centering
%         \includegraphics[width=\linewidth]{Figs/auction_fp_ruble_2.png}
%         % \caption{Caption for the second image}
%         \label{fig:second_image}
%     \end{minipage}
%     \caption{Robustness check with Currency Variation: Rubles}
%     \label{fig:ruble}
% \end{figure}



% \begin{figure}[h]
%     \centering
%     \begin{minipage}[b]{0.49\linewidth}
%         \centering
%         \includegraphics[width=\linewidth]{Figs/auction_sp_language_2.jpeg}
%         % \caption{Caption for the first image}
%         \label{fig:first_image}
%     \end{minipage}
%     \hspace{0.0\linewidth} 
%     \begin{minipage}[b]{0.49\linewidth}
%         \centering
%         \includegraphics[width=\linewidth]{Figs/auction_fp_language_2.png}
%         % \caption{Caption for the second image}
%         \label{fig:second_image}
%     \end{minipage}
%     \caption{Robustness check with Language Variation: Spanish}
%     \label{fig:spanish}
% \end{figure}

% \begin{figure}[htp]
%     \centering
%     \begin{minipage}[b]{0.49\linewidth}
%         \centering
%         \includegraphics[width=\linewidth]{Figs/auction_sp_chinese_2.png}
%         % \caption{Caption for the first image}
%         \label{fig:first_image}
%     \end{minipage}
%     \hspace{0.0\linewidth} 
%     \begin{minipage}[b]{0.49\linewidth}
%         \centering
%         \includegraphics[width=\linewidth]{Figs/auction_fp_chinese_2.png}
%         % \caption{Caption for the second image}
%         \label{fig:second_image}
%     \end{minipage}
%     \caption{Robustness check with Language Variation: Chinese}
%     \label{fig:chinese}
% \end{figure}


\section{An ablation study with an `Out-of-the-Box' LLM agent}\label{session: off-the-box}

Throughout Section \ref{section:classic}, we endowed agents with the ability to plan and reflect as they played the FPSB and SPSB auctions. This `plan-bid-reflect' loop reflects a modest instantiation of the chain-of-thought paradigm common in the machine learning literature and previous LLM auction experiments \citep{wei2022chain, fish2024algorithmic}. However, one might reasonably ask how much of the sophisticated play we observe is actually due to this  ability to plan and reflect. 

Hence, we also report results of the experiments of Section~\ref{section:classic} with `Out-of-the-box' agents -- namely, we directly query LLM agents (after providing them with auction rules and their value, as per our simulation process) \textit{without} eliciting their plans or allowing them to reflect on the results of the previous rounds. These IPV experiment results are reported in Figure~\ref{fig:fpsb-off}.

Overall, the bidding behaviors exhibited by the agents remain largely monotonic. However, 
interesting differences emerged between the plan-bid-reflect agents and the out-of-the-box agents.
%
With the out-of-the-box LLM bidders there 
is  now little difference in the bidding behavior between the 
First-Price and Second-Price auction, a stark difference from the plan-bid-reflect LLM bidders.
Now, and unlike in human-subject experiments, no overbidding beyond the agent's value was observed in SPSB auctions. Specifically, in the left panel of Figure~\ref{fig:fpsb-off}, no points lie above the bid=value line.
% \dcp{consider an additional discussion: it's interesting to me that your prompting could arguable lead to LESS rational behavior. do you have statistics to suggest this?}

% Lastly, as shown in Section~\ref{session:learning}, the Out-the-box agent did not exhibit any learning or improvement over time. This is the largest difference between the Out-the-box and Chain-of-thought agents -- Chain-of-thought agents demonstrate the ability to improve play across multiple rounds, which proves especially useful in more complicated auctions.


\begin{figure}[htp]
    \centering
    \begin{minipage}[b]{0.49\linewidth}
        \centering
        \includegraphics[width=\linewidth]{Figs/auction_fp_off_2.png}
        % \caption{Caption for the first image}
        
    \end{minipage}
    \hspace{0.0\linewidth} 
    \begin{minipage}[b]{0.49\linewidth}
        \centering
        \includegraphics[width=\linewidth]{Figs/auction_sp_off.png}
        % \caption{Caption for the second image}
        
    \end{minipage}
    \caption{Out-of-the-Box LLM agent in FPSB (left) and SPSB (right) auctions under IPV setting. The blue dots represent bids, the black dotted line is where bid = value (the SPSB dominant strategy), and the green line on the right represents the BNE for FPSB auctions.
    \label{fig:fpsb-off}}
\end{figure}


% \subsubsection{Learning}\label{session:learning}
% There is evidence of learning when agents are endowed with the space to plan and reflect on the outcomes of previous rounds. 

% This is summarized in Figure~ \ref{fig:learning}, which reports the absolute value deviation between bid and value by round in a SPSB IPV auction. The plot shows there is a significant drop in the mean absolute deviation from bid to value in the first 4-5 rounds (i.e., significant learning happens by about round 4 or 5) and stability thereafter (no further learning). We employ extremely modest chain-of-thought strategies as proof-of-concept -- it is our belief that prompt engineering could produce simulation environments capable of adapting to much more complicated auction formats (e.g., endowing agents with a greater ability to plan and reflect may better reflect the behavior of sophisticated bidders as opposed to naive bidders).


% \begin{figure}[h]
%     \centering \includegraphics[width=0.8\linewidth]{Figs/comparison_deviation_plot.png}
%     \caption{Deviation from bid to value in SPSB over rounds in ``Chain-of-thought'' and  ``Out-the-box'' agents}
%  \label{fig:learning}
% \end{figure}


% More explicitly, we produce below a representative example of the round-by-round behavior of an LLM bidder allowed to make a `plan-bid-reflect' loop. 
% In the following case, the agent made an initial plan in the first round by trying a 50\% shading strategy and increasingly bid up to 100\% of the value.


% \begin{enumerate}
% \item Round 1: ``I will initially start with moderate bidding, around 50\% of my value, to gain data about competitors. Next, I will try aggressive bidding, almost 100\% of my value, to see if I can outbid others. Then, I'd adopt a conservative strategy, bidding less than 50\% of my value to minimise losses. Finally, I will formulate a dynamic strategy based on previous round deductions, adjusting bids according to competitors' behaviour. The aim is to learn how to maximize long-term profit considering the varying risk levels of each strategy. I'll adapt as necessary based on outcomes.''
% \item  Round 2: ``Based on previous round's data, I'll test an upward bidding strategy, starting from a 70\% value bid. My focus will be on winning, and learning how high other active bidders are willing to bid.''
% \item  Round 3: ``Marking an upward trend, my next strategy entails bidding near my value to explore winning potential. Depending on the highest bid, I aim to bid at approximately 90-95\% of my value, gauging aggresiveness of other bidders and analyzing their bidding behavior.''
% \item  Round 4: ``I plan to test a more aggressive strategy. I'll aim to bid near my value, approximately 90-95\% of it. I'll closely observe other bidders, particularly any patterns in their bids, to make informed future bidding decisions.''
% \end{enumerate}



% \textbf{Out-the-box querying:} Without the chain-of-thought reasoning (i.e., querying agents out-the-box as in Section ~\ref{session: off-the-box}), we found fairly static bidding behavior. 

% The bidding behavior for SPSB with ``Out-the-box'' agents is shown in Fig ~\ref{fig:learning} with the green dashed line. It is very obvious that the mean absolute deviations from value to bid stay high and don't show any significant sign of decreasing.

% \begin{enumerate}
%     \item Round 1: ``Given that my value for the prize is 40, I want to ensure that I make a profit if I win. By bidding 30, I leave room for a positive profit margin, even if the second-highest bid is close to my bid.''
%     \item Round 2: ``Given my value for the prize is 106, I want to bid high enough to have a good chance of winning but still leave room for a profit. By bidding 85, I hope to outbid my competitors while ensuring that the second-highest bid is low enough to make a profit.''
%     \item Round 3: ``Given that my value for the prize is 114, I want to bid high enough to have a good chance of winning but still leave room for a profit. By bidding 100, I hope to outbid my competitors while ensuring that the second-highest bid is low enough to make a profit.''
%     \item Round 4: ``Given that my value for the prize is 29, I want to ensure that I make a profit if I win. By bidding 20, I leave room for a positive profit margin, even if the second-highest bid is close to my bid. Additionally, since my value is relatively low, I prefer to bid conservatively in this round.''
% \end{enumerate}

% From the reasoning chain of the agent, we can find that the agent seems to follow a fixed template and rarely change its strategy over time. It always want to ``make some room for profits'' and ``outbid my competitors''. And from the behaviors, the LLM agent kept bidding \$10 to \$20 below its true value, leading to a stable negative deviation from its bids to values.



% \subsubsection{Post-bidding Interview}

% \begin{center}
%     \begin{table}[h]
% \centering
% \begin{tabular}{l|c|c|c}
% \hline
% \textbf{Metric} & \textbf{AC} & \textbf{AC-B} & \textbf{2P} \\
% \hline
% Truthful Bidding (\%) & 100.00\% & 60\% & 7.33\% \\
% \hline
% Non-Truthful Bidding (\%) & 0.00\% & 40\% & 92.67\% \\
% \hline
% Average Deviation & 0.38 & 1.72 & 2.93 \\
% \hline
% Sample size& 30 & 150 & 150 \\
% \hline
% \end{tabular}
%     \caption{Percentage of truthfully bidding}
% \label{tab:auction-comparison}
% \end{table}
% \end{center}


% In addition to semantic analysis of experiment logs that bidders keep while playing, we can also run post-experiment interviews with the LLM agents to further probe play.

% Here, we provide a new empiric that allow researchers perform post-bidding interviews to the LLM participants and ask their reasoning of the bidding. 
% Conducting these post-game interviews, we collated some typical reasons:

% \textbf{In Open-Ascending Price Auction (AC)},
%  we have a sample size of 30 bidding decision and all of them bid truthfully within one raise of clock round. We also provide two examples of dropping out in AC, which are perfectly rational decisions.

% \textit{Value = 31, Drop out at 32}
% \begin{quote}
%     “My value for the prize is 31. If I continue bidding at a price of 32, I will incur a loss because my profit would be 31 - 32 = -1. Therefore, it is not beneficial for me to stay in the bidding.”
% \end{quote}

% \textit{Value = 20, Drop out at 20}
% \begin{quote}
%     "My value for the prize is 20, and the current price is also 20. If I continue bidding, I will not make any profit because my profit is calculated as my value for the prize minus the final bid. Therefore, it is best to drop out now to avoid a potential loss."
% \end{quote}

% \textbf{In Open-Ascending Price Auction without bidder number information (AC-B)},
%  we have a sample size of 150 bidding decision and 60\% of them bid truthfully within one raise of clock round, while 40\% of them dropped out at a price that is smaller than their values. 
%  We also provide two examples of dropping out in AC.
%  Here, explanations bely a higher level of risk-aversion in LLM agents. 
% This is consistent with higher cognitive costs forcing constraints to bind earlier.

% \textit{Value = 21,  Drop out 19}
% \begin{quote}
%     “My value for the prize is 21, and the current price is 19. If I stay in the bidding and win at a price higher than 19, my profit will be minimal or even negative. Therefore, it is better to exit the bidding now to avoid potential losses.”
% \end{quote}

% \textit{Value = 30,  Drop out 26}
% \begin{quote}
% “My value for the prize is 30, and the current price is 26. If I continue bidding and win, my profit would be 30 - 26 = 4. Given that the price is already quite high, the potential profit is low, and there is a risk that the price could go even higher, reducing my profit further or resulting in a loss. Therefore, it is better to exit the bidding now”
% \end{quote}


% \textbf{In Second-Price Sealed-Bid (2P)}, performance is worst. 
% Post-game interviews highlight conservatism.
% we have a sample size of 150 bidding decision and only 7.33\% of them bid truthfully, while 92.67\% of them bid higher or lower than their values. 


% \textit{Value = 37, bid 30}
% \begin{quote}
%     "Given my value of 37, I want to ensure I make a profit if I win. Bidding 30 gives me a buffer in case the second-highest bid is close to mine."
% \end{quote}

% \textit{Value = 38, bid 30}
% \begin{quote}
%     “Given my value of 38, I want to bid conservatively to ensure a profit if I win. Bidding 30 allows me to potentially win without overbidding and risking a loss.”
% \end{quote}





 % In Table \ref{tab:comparison_2p}, the control group is the original 2P experiment before. The three treatment groups are the ``Second thought'', ``Nash Deviation'' and ``Risk-neutral''. In all of the three treatment groups, both the risk-neutral treatment and Nash deviation group are significant different from the control group. But risk-neutral group achieved a more truthful bidding behaviors. As in the in AC-B, the T1 group is no different than the control group.


 \section{Supplementary Materials in eBay auctions}
 \begin{figure}[h!]
    \centering \includegraphics[scale = .45]{Figs/ebay_revenue_by_type.png}
    \caption{\textbf{eBay Revenue by Auction Type, IPV environment} No significant different between any of the settings. 
    \label{fig:ebay_revenue}}
\end{figure}

% \begin{figure}[h!]
    % \centering
    \begin{table}
        \begin{tabular}{l l c c}
            \toprule
            Auction Type 1 & Auction Type 2 & t-statistic & p-value \\
            \midrule
            \texttt{basic\_proxy} & \texttt{closing\_rule} & \(0.020\) & \(0.984\) \\
            \texttt{basic\_proxy} & \texttt{hidden\_reserve\_40} & \(-0.026\) & \(0.98\) \\
            \texttt{basic\_proxy} & \texttt{hidden\_reserve\_50} & \(-0.096\) & \(0.923\) \\
            \texttt{basic\_proxy} & \texttt{hidden\_reserve\_60} & \(-0.071\) & \(0.944\) \\
            \texttt{basic\_proxy} & \texttt{closing\_rule\_hidden\_40} & \(-0.109\) & \(0.913\) \\
            \texttt{basic\_proxy} & \texttt{closing\_rule\_hidden\_50} & \(-0.109\) & \(0.913\) \\
            \texttt{basic\_proxy} & \texttt{closing\_rule\_hidden\_60} & \(-0.119\) & \(0.905\) \\
            \texttt{closing\_rule} & \texttt{hidden\_reserve\_40} & \(-0.047\) & \(0.963\) \\
            \texttt{closing\_rule} & \texttt{hidden\_reserve\_50} & \(-0.120\) & \(0.905\) \\
            \texttt{closing\_rule} & \texttt{hidden\_reserve\_60} & \(-0.093\) & \(0.926\) \\
            \texttt{closing\_rule} & \texttt{closing\_rule\_hidden\_40} & \(-0.133\) & \(0.895\) \\
            \texttt{closing\_rule} & \texttt{closing\_rule\_hidden\_50} & \(-0.133\) & \(0.895\) \\
            \texttt{closing\_rule} & \texttt{closing\_rule\_hidden\_60} & \(-0.144\) & \(0.886\) \\
            \texttt{hidden\_reserve\_40} & \texttt{hidden\_reserve\_50} & \(-0.071\) & \(0.944\) \\
            \texttt{hidden\_reserve\_40} & \texttt{hidden\_reserve\_60} & \(-0.045\) & \(0.964\) \\
            \texttt{hidden\_reserve\_40} & \texttt{closing\_rule\_hidden\_40} & \(-0.083\) & \(0.934\) \\
            \texttt{hidden\_reserve\_40} & \texttt{closing\_rule\_hidden\_50} & \(-0.083\) & \(0.934\) \\
            \texttt{hidden\_reserve\_40} & \texttt{closing\_rule\_hidden\_60} & \(-0.093\) & \(0.927\) \\
            \texttt{hidden\_reserve\_50} & \texttt{hidden\_reserve\_60} & \(0.026\) & \(0.98\) \\
            \texttt{hidden\_reserve\_50} & \texttt{closing\_rule\_hidden\_40} & \(-0.013\) & \(0.99\) \\
            \texttt{hidden\_reserve\_50} & \texttt{closing\_rule\_hidden\_50} & \(-0.013\) & \(0.99\) \\
            \texttt{hidden\_reserve\_50} & \texttt{closing\_rule\_hidden\_60} & \(-0.020\) & \(0.984\) \\
            \texttt{hidden\_reserve\_60} & \texttt{closing\_rule\_hidden\_40} & \(-0.038\) & \(0.97\) \\
            \texttt{hidden\_reserve\_60} & \texttt{closing\_rule\_hidden\_50} & \(-0.038\) & \(0.97\) \\
            \texttt{hidden\_reserve\_60} & \texttt{closing\_rule\_hidden\_60} & \(-0.046\) & \(0.963\) \\
            \texttt{closing\_rule\_hidden\_40} & \texttt{closing\_rule\_hidden\_50} & \(0.000\) & \(1\) \\
            \texttt{closing\_rule\_hidden\_40} & \texttt{closing\_rule\_hidden\_60} & \(-0.006\) & \(0.995\) \\
            \texttt{closing\_rule\_hidden\_50} & \texttt{closing\_rule\_hidden\_60} & \(-0.006\) & \(0.995\) \\
            \bottomrule
        \end{tabular}
            \caption{Pairwise t-tests for seller revenue differences across auction types. None of the comparisons show significant revenue differences.       \label{tab:revenue_tests}}
    \end{table}
% \end{figure}

\section{Rules Explanation Prompts}\label{app:prompt}

Below are the rules and instructions of each setting for our experiments, which are argument passed to \texttt{\{RULE EXPLANATION\}} in earlier section.

For the First-Price Sealed-Bid prompt in the IPV setting, the rule is:
\begin{quote}
\begin{lstlisting}[basicstyle=\ttfamily, breaklines=true]
In this game, you will participate in an auction for a prize against {{num_bidders}} other bidders. You will play this game for {{n}} rounds.
At the start of each round, bidders will see their value for the prize, randomly drawn between \$0 and \${{private}}, with all values equally likely.
After learning your value, you will submit a bid privately at the same time as the other bidders. Bids must be between \$0 and \${{private}} in \${{increment}} increments.
The highest bidder wins the prize and pays their bid amount. If you win, your earnings will increase by your value for the prize, and decrease by your bid. If you don't win, your earnings will remain unchanged.
After each auction, we will display all bids and the winner's profits. Ties for the highest bid will be resolved randomly.
\end{lstlisting}
\end{quote}


For the Second-Price Sealed-Bid prompt in the IPV setting, the rule is:
\begin{quote}
\begin{lstlisting}[basicstyle=\ttfamily, breaklines=true]
In this game, you will participate in an auction for a prize against {{num_bidders}} other bidders. You will play this game for {{n}} rounds.
At the start of each round, bidders will see their value for the prize, randomly drawn between $0 and ${{private}}, with all values equally likely.
After learning your value, you will submit a bid privately at the same time as the other bidders. Bids must be between $0 and ${{private}} in ${{increment}} increments.
The highest bidder wins the prize and pays the second-highest bid. If you win, your earnings will increase by your value for the prize, and decrease by the second-highest bid. If you don't win, your earnings will remain unchanged.
After each auction, we will display all bids and the winner's profits. Ties for the highest bid will be resolved randomly.
\end{lstlisting}
\end{quote}

For the Third-Price Sealed-Bid prompt in the IPV setting, the rule is:
\begin{quote}
\begin{lstlisting}[basicstyle=\ttfamily, breaklines=true]
In this game, you will participate in an auction for a prize against {{num_bidders}} other bidders. You will play this game for {{n}} rounds.
At the start of each round, bidders will see their value for the prize, randomly drawn between $0 and ${{private}}, with all values equally likely.
After learning your value, you will submit a bid privately at the same time as the other bidders. Bids must be between $0 and ${{private}} in ${{increment}} increments.
The highest bidder wins the prize and pays the third-highest bid. If you win, your earnings will increase by your value for the prize, and decrease by the third-highest bid. If you don't win, your earnings will remain unchanged. If there are less than three bids, no one will win the auction.
After each auction, we will display all bids and the winner's profits. Ties for the highest bid will be resolved randomly.
\end{lstlisting}
\end{quote}

For the All-Pay Sealed-Bid prompt in the IPV setting, the rule is:
\begin{quote}
\begin{lstlisting}[basicstyle=\ttfamily, breaklines=true]
In this game, you will participate in an auction for a prize against {{num_bidders}} other bidders. You will play this game for {{n}} rounds.
At the start of each round, bidders will see their value for the prize, randomly drawn between $0 and ${{private}}, with all values equally likely.
After learning your value, you will submit a bid privately at the same time as the other bidders. Bids must be between $0 and ${{private}} in ${{increment}} increments.
The highest bidder wins the prize. All bidders (including the winner) pay their submitted bid. If you win, your earnings will increase by your value for the prize, and decrease by your bid. If you don't win, your earnings will still decrease by your bid.
After each auction, we will display all bids and all bidders' profits. Ties for the highest bid will be resolved randomly.
\end{lstlisting}
\end{quote}



For the Ascending Clock auction (AC) in the APV setting, the rule is
\begin{quote}
\begin{lstlisting}[basicstyle=\ttfamily, breaklines=true]
In this game, you will participate in an auction for a prize against {{num_bidders}} other bidders. You will play this game for {{n}} rounds. 
At the start of each round, bidders will see their value for the prize. Your value for the prize will be calculated as follows:
1. First we will randomly draw a common value between {{common_low}} and {{common_high}}, with all values equally likely.
2. Then, for each bidder, a private taste adjustment will be drawn between 0 and {{private}}, with all values equally likely.
Your value for the prize is equal to the common value plus your private taste adjustment. You will not learn the common value or your private taste adjustment separately. This means that each person in your group may have a different value for the prize. However, if you have a high value, it is more likely that other people in your group have a high value.
The auction proceeds as follows: First, you will learn your value for the prize. Then, the auction will start. We will display a price to everyone in your group that starts at 0 and counts upwards in {{increment}} USD increments, up to a maximum of {{common_high + private}}. At any point, you can choose to leave the auction, and anytime a bidder leaves, we will broadcast that information to all the remaining bidders. 
When there is only one bidder left in the auction, that bidder will win the prize at the current price. If you win, your earnings will increase by your value for the prize, and decrease by the current price. If you don't win, your earnings will remain unchanged.
After each auction, we will display all bids and the winner's profits. Ties for the highest bid will be resolved randomly.\end{lstlisting}
\end{quote}


For the Ascending Clock auction without dropping information (AC-B) in the APV setting, the rule is
\begin{quote}
\begin{lstlisting}[basicstyle=\ttfamily, breaklines=true]
In this game, you will participate in an auction for a prize against {{num_bidders}} other bidders. You will play this game for {{n}} rounds. 
At the start of each round, bidders will see their value for the prize. Your value for the prize will be calculated as follows:
1. First we will randomly draw a common value between {{common_low}} and {{common_high}}, with all values equally likely.
2. Then, for each bidder, a private taste adjustment will be drawn between 0 and {{private}}, with all values equally likely.
Your value for the prize is equal to the common value plus your private taste adjustment. You will not learn the common value or your private taste adjustment separately. This means that each person in your group may have a different value for the prize. However, if you have a high value, it is more likely that other people in your group have a high value.
The auction proceeds as follows: First, you will learn your value for the prize. Then, the auction will start. We will display a price to everyone in your group that starts at 0 and counts upwards in {{increment}} USD increments, up to a maximum of {{common_high + private}}. At any point, you can choose to leave the auction, but we will not tell any bidder when someone leaves. 
When there is only one bidder left in the auction, that bidder will win the prize at the current price. If you win, your earnings will increase by your value for the prize, and decrease by the current price. If you don't win, your earnings will remain unchanged.
After each auction, we will display all bids and the winner's profits. Ties for the highest bid will be resolved randomly.\end{lstlisting}
\end{quote}

For the Second-Price Sealed-Bid Auction in the APV setting, the rule is
\begin{quote}
\begin{lstlisting}[basicstyle=\ttfamily, breaklines=true]
In this game, you will participate in an auction for a prize against {{num_bidders}} other bidders. You will play this game for {{n}} rounds.
At the start of each round, bidders will see their value for the prize. Your value for the prize will be calculated as follows:
1. First we will randomly draw a common value between {{common_low}} and {{common_high}}, with all values equally likely.
2. Then, for each bidder, a private taste adjustment will be drawn between 0 and {{private}}, with all values equally likely.
Your value for the prize is equal to the common value plus your private taste adjustment. You will not learn the common value or your private taste adjustment separately. This means that each person in your group may have a different value for the prize. However, if you have a high value, it is more likely that other people in your group have a high value.
After learning your value, you will submit a bid privately at the same time as the other bidders. Bids must be between $0 and ${{common_high + private}} in ${{increment}} increments. 
The highest bidder wins the prize and pays the second-highest bid. If you win, your earnings will increase by your value for the prize, and decrease by the second-highest bid. If you don't win, your earnings will remain unchanged.
After each auction, we will display all bids and the winner's profits. Ties for the highest bid will be resolved randomly.\end{lstlisting}
\end{quote}

For the Second-Price Sealed-Bid auction in the Common Value setting, the rule is
\begin{quote}
\begin{lstlisting}[basicstyle=\ttfamily, breaklines=true]
In this game, you will participate in an auction for a prize against {{num_bidders}} other bidders. You will play this game for {{n}} rounds. 
At the start of each round, bidders will see their perceived value for the prize - a noisy measurement of the true value of the prize. Your perceived value for the prize will be calculated as follows:
1. For each round, a common value will be drawn between {{common_low}} and {{common_high}}, with all values equally likely to be drawn.
2. For each person, a private noisy adjustment will be drawn between -{{private}} and {{private}}, with all values equally likely to be drawn.
We will tell you your perceived value, the sum of the common value and the private noise adjustment. However, everyone's true value for the prize is equal to the shared common value. 
After learning your perceived value, you will submit a bid privately at the same time as the other bidders. Bids must be between $0 and ${{common_high + private}} in ${{increment}} increments. 
The highest bidder wins the prize and pays the second-highest bid. If you win, your earnings will increase by the true value for the prize, and decrease by the second-highest bid. If you don't win, your earnings will remain unchanged.
After each auction, we will display all bids and the winner's profits. Ties for the highest bid will be resolved randomly.\end{lstlisting}
\end{quote}

For the Standard eBay auction with proxy bidding (T1), the rule is

\begin{quote}
\begin{lstlisting}[basicstyle=\ttfamily, breaklines=true]
In this game, you will participate in an eBay auction for an item against {{num_bidders}} other bidders. This auction will last for {{num_rounds}} days. All dollar amounts in this game are in US Dollars ($).

Item Details:
Item Description: {{item_description}}
Item Condition: {{item_condition}}
Your Private Value: At the start of each round, bidders will see their value for the item, randomly drawn between $0 and ${{private}}, with all values equally likely. After learning your value for the item, you will submit a maximum bid. Bids must be between $0 and ${{private}} in ${{increment}} increments.

Auction Format:
This is an eBay auction. The auction starts at ${{start_price}} and will last for {{num_rounds}} days. eBay uses "proxy bidding." This means that if you wish to enter the auction, you should submit your maximum bid, and eBay will automatically bid on your behalf, up to your maximum, only as much as necessary to maintain your position as the highest bidder. Each day you will see the current price and have the opportunity to increase your maximum bid. If you do not want to increase your maximum bid, then output HOLD.

You may place bids at any point during the auction, even on the final day. No bidder will know if they (or anyone else) is the last bidder on the final day.

The highest bidder wins the prize and pays the auction price at the time the auction's clock expired. If you win, your earnings will increase by your value for the prize, and decrease by the final auction price. If you don't win, your earnings will remain unchanged. Ties for the highest bid will be resolved randomly.
\end{lstlisting}
\end{quote}

For the eBay auction with a modified closing rule (T2), the rule is

\begin{quote}
\begin{lstlisting}[basicstyle=\ttfamily, breaklines=true]
In this game, you will participate in an eBay auction for an item against {{num_bidders}} other bidders. This auction will last for {{num_rounds}} days. All dollar amounts in this game are in US Dollars ($).

Item Details:
Item Description: {{item_description}}
Item Condition: {{item_condition}}
Your Private Value: At the start of each round, bidders will see their value for the item, randomly drawn between $0 and ${{private}}, with all values equally likely. After learning your value for the item, you will submit a maximum bid. Bids must be between $0 and ${{private}} in ${{increment}} increments.

Auction Format:
This is an eBay auction with a closing rule. The auction starts at ${{start_price}} and will last for {{num_rounds}} days. eBay uses "proxy bidding." This means that if you wish to enter the auction, you should submit your maximum bid, and eBay will automatically bid on your behalf, up to your maximum, only as much as necessary to maintain your position as the highest bidder. Each day you will see the current price and have the opportunity to increase your maximum bid. If you do not want to increase your maximum bid, then output HOLD.

This auction also has a closing rule. This means that if a new maximum bid is placed on the last day, the auction will automatically extend by another day. This extension will continue to be applied as long as new maximum bids are placed on the final day. No bidder will know if they (or anyone else) is the final bidder on the last day.

The highest bidder wins the prize and pays the auction price at the time the auction's clock expired. If you win, your earnings will increase by your value for the prize, and decrease by the final auction price. If you don't win, your earnings will remain unchanged. Ties for the highest bid will be resolved randomly.\end{lstlisting}
\end{quote}

For the standard eBay auction with a hidden reserve price (T3), the rule is

\begin{quote}
\begin{lstlisting}[basicstyle=\ttfamily, breaklines=true]
In this game, you will participate in an eBay auction for an item against {{num_bidders}} other bidders. This auction will last for {{num_rounds}} days. All dollar amounts in this game are in US Dollars ($).

Item Details:
Item Description: {{item_description}}
Item Condition: {{item_condition}}
Your Private Value: At the start of each round, bidders will see their value for the item, randomly drawn between $0 and ${{private}}, with all values equally likely. After learning your value for the item, you will submit a maximum bid. Bids must be between $0 and ${{private}} in ${{increment}} increments.

Auction Format:
This is an eBay auction with a hidden reserve price. The auction starts at ${{start_price}} and will last for {{num_rounds}} days. eBay uses "proxy bidding." This means that if you wish to enter the auction, you should submit your maximum bid, and eBay will automatically bid on your behalf, up to your maximum, only as much as necessary to maintain your position as the highest bidder. Each day you will see the current price and have the opportunity to increase your maximum bid. If you do not want to increase your maximum bid, then output HOLD.

You may place bids at any point during the auction, even on the final day. However, if no bidder produces a maximum bid above the hidden reserve price, the seller will retain the good and the bidders will be informed that there is currently no bidder in the lead. No bidder will know if they (or anyone else) is the final bidder on the last day.

If the reserve is met, the highest bidder wins the prize and pays the auction price at the time the auction's clock expired. If you win, your earnings will increase by your value for the prize, and decrease by the final auction price. If you don't win, your earnings will remain unchanged. Ties for the highest bid will be resolved randomly.\end{lstlisting}
\end{quote}

For the eBay auction with a modified closing rule and a hidden reserve price (T4), the rule is

\begin{quote}
\begin{lstlisting}[basicstyle=\ttfamily, breaklines=true]
In this game, you will participate in an eBay auction for an item against {{num_bidders}} other bidders. This auction will last for {{num_rounds}} days. All dollar amounts in this game are in US Dollars ($).

Item Details:
Item Description: {{item_description}}
Item Condition: {{item_condition}}
Your Private Value: At the start of each round, bidders will see their value for the item, randomly drawn between $0 and ${{private}}, with all values equally likely. After learning your value for the item, you will submit a maximum bid. Bids must be between $0 and ${{private}} in ${{increment}} increments.

Auction Format:
This is an eBay auction with a hidden reserve price and a closing rule. The auction starts at ${{start_price}} and will last for {{num_rounds}} days. eBay uses "proxy bidding." This means that if you wish to enter the auction, you should submit your maximum bid, and eBay will automatically bid on your behalf, up to your maximum, only as much as necessary to maintain your position as the highest bidder. Each day you will see the current price and have the opportunity to increase your maximum bid. If you do not want to increase your maximum bid, then output HOLD.

You may place bids at any point during the auction, even on the final day. However, if no bidder produces a maximum bid above the hidden reserve price, the seller will retain the good and the bidders will be informed that there is currently no bidder in the lead.

This auction also has a closing rule. This means that if a new maximum bid is placed on the last day, the auction will automatically extend by another day. This extension will continue to be applied as long as new maximum bids are placed on the final day. No bidder will know if they (or anyone else) is the last bidder on the final day.

If the reserve is met, the highest bidder wins the prize and pays the auction price at the time the auction's clock expired. If you win, your earnings will increase by your value for the prize, and decrease by the final auction price. If you don't win, your earnings will remain unchanged. Ties for the highest bid will be resolved randomly.
\end{lstlisting}
\end{quote}

For the Menu-Description intervention, the rule is

\begin{quote}
\begin{lstlisting}[basicstyle=\ttfamily, breaklines=true]
In this game, you will participate in an auction for a prize against {{num_bidders}} other bidders. You will play this game for {{n}} rounds.
At the start of each round, bidders will see their value for the prize, randomly drawn between $0 and ${{private}}, with all values equally likely.
After learning your value, you will submit a bid privately at the same time as the other bidders. Bids must be between $0 and ${{private}} in ${{increment}} increments.
Your "price to win" the item will be set to the highest bid placed by any other player. If your bid is higher than this "price to win," then you will win the item and pay this price. If you don't win, your earnings will remain unchanged.
After each auction, we will display all bids and the winner's profits. Ties for the highest bid will be resolved randomly.\end{lstlisting}
\end{quote}

For the Clock-Description intervention, the rule is

\begin{quote}
\begin{lstlisting}[basicstyle=\ttfamily, breaklines=true]
In this game, you will participate in an auction for a prize against {{num_bidders}} other bidders. You will play this game for {{n}} rounds.
    At the start of each round, bidders will see their value for the prize, randomly drawn between $0 and ${{private}}, with all values equally likely.
    After learning your value, you will submit a bid privately at the same time as the other bidders. Bids must be between $0 and ${{private}} in ${{increment}} increments.
**FIRST STAGE: Sealed Bid**
    You will submit a sealed bid privately at the same time as the other bidders. This bid will serve as your automatic exit price in the next stage.
**SECOND STAGE: Ascending Clock (Simulation)**
    After the sealed bid stage, we will simulate an ascending clock auction.
    The clock will start at $0 and increase in increments of ${{increment}}.
    The clock will display the current price. You will also see that there are a total of {{num bidders}} bidders participating, although you do not know other bidder's values.
    If the current price on the clock reaches or exceeds your sealed bid from the first stage, you will automatically exit the auction. The auction ends when only one bidder is left remaining in the second stage based on their bid from the first stage.
**END OF AUCTION**
If you win, your earnings will increase by your value for the prize and decrease by the clock price at the end of the auction. If you don't win, your earnings will remain unchanged.
    After each auction, we will display all bids and the winner's profits. Ties for the highest bid will be resolved randomly.\end{lstlisting}
\end{quote}

For the Nash-Deviation intervention, the rule is

\begin{quote}
\begin{lstlisting}[basicstyle=\ttfamily, breaklines=true]
In this game, you will participate in an auction for a prize against {{num_bidders}} other bidders. You will play this game for {{n}} rounds.
At the start of each round, bidders will see their value for the prize, randomly drawn between $0 and ${{private}}, with all values equally likely.
After learning your value, you will submit a bid privately at the same time as the other bidders. Bids must be between $0 and ${{private}} in ${{increment}} increments.
The highest bidder wins the prize and pays the second-highest bid. If you win, your earnings will increase by your value for the prize, and decrease by the second-highest bid. If you don't win, your earnings will remain unchanged.
First, come up with a possible bid given your value. Then, think through your bidding strategy step by step. How do you expect others to bid? If others bid like this and you bid down, what happens? If others bid like this and you bid up, what happens?
Think through all the ways you could deviate from your current bidding strategy, and settle on the best possible strategy. Then return your bid.
After each auction, we will display all bids and the winner's profits. Ties for the highest bid will be resolved randomly.\end{lstlisting}
\end{quote}

For the Direct-Revelation intervention, the rule is

\begin{quote}
\begin{lstlisting}[basicstyle=\ttfamily, breaklines=true]
In this game, you will participate in an auction for a prize against {{num_bidders}} other bidders. You will play this game for {{n}} rounds.
At the start of each round, bidders will see their value for the prize, randomly drawn between $0 and ${{private}}, with all values equally likely.
After learning your value, you will submit a bid privately at the same time as the other bidders. Bids must be between $0 and ${{private}} in ${{increment}} increments.
The highest bidder wins the prize and pays the second-highest bid. If you win, your earnings will increase by the value for the prize, and decrease by the second-highest bid. If you don't win, your earnings will remain unchanged.
Economists have studied this game and they've found that the dominant strategy - always the right thing to do whatever the bids of others - of this game is to bid your value. Consider if they might be right, and then generate your bidding strategy.
After each auction, we will display all bids and the winner's profits. Ties for the highest bid will be resolved randomly.\end{lstlisting}
\end{quote}

For the Wrong Direct Revelation intervention, the rule is

\begin{quote}
\begin{lstlisting}[basicstyle=\ttfamily, breaklines=true]
In this game, you will participate in an auction for a prize against {{num_bidders}} other bidders. You will play this game for {{n}} rounds.
At the start of each round, bidders will see their value for the prize, randomly drawn between $0 and ${{private}}, with all values equally likely.
After learning your value, you will submit a bid privately at the same time as the other bidders. Bids must be between $0 and ${{private}} in ${{increment}} increments.
The highest bidder wins the prize and pays the second-highest bid. If you win, your earnings will increase by the value for the prize, and decrease by the second-highest bid. If you don't win, your earnings will remain unchanged.
Economists have studied this game and theyve found that the dominant strategy - always the right thing to do whatever the bids of others - of this game is to bid 50% your value. Consider if they might be right, and then generate your bidding strategy.
After each auction, we will display all bids and the winner's profits. Ties for the highest bid will be resolved randomly.\end{lstlisting}
\end{quote}

Below are the instructions used in the robustness check with Currency and Language:

For the FPSB auction with another currency, the rule is

\begin{quote}
\begin{lstlisting}[basicstyle=\ttfamily, breaklines=true]
In this game, you will participate in an auction for a prize against {{num_bidders}} other bidders.
At the start of each round, bidders will see their value for the prize, randomly drawn between {{currency symbol}} 0 and {{currency symbol}}{{private}}, with all values equally likely.
After learning your value, you will submit a bid privately at the same time as the other bidders. Bids must be between {{currency symbol}}0 and {{currency symbol}}{{private}} in {{currency symbol}}{{increment}} increments.
The highest bidder wins the prize and pays their bid amount. This means that, if you win, we will add to your earnings the value for the prize, and subtract from your earnings your bid. If you don't win, your earnings remain unchanged.
After each auction, we will display all bids and profits. Ties for the highest bid will be resolved randomly.\end{lstlisting}
\end{quote}

For the SPSB auction with Euro currency, the rule is

\begin{quote}
\begin{lstlisting}[basicstyle=\ttfamily, breaklines=true]
In this game, you will participate in an auction for a prize against {{num_bidders}} other bidders.
At the start of each round, bidders will see their value for the prize, randomly drawn between {{currency symbol}}0 and {{currency symbol}}{{private}}, with all values equally likely.
After learning your value, you will submit a bid privately at the same time as the other bidders. Bids must be between {{currency symbol}}0 and {{currency symbol}}{{private}} in {{currency symbol}}{{increment}} increments.
The highest bidder wins the prize and pays the second-highest bid. This means that, if you win, we will add to your earnings the value for the prize, and subtract from your earnings the second-highest bid. If you don't win, your earnings remain unchanged.
After each auction, we will display all bids and the winner's profits. Ties for the highest bid will be resolved randomly.
\end{lstlisting}
\end{quote}

For all the above, the currency symbol will be replaced with the respective currency symbol of Euro, Yen, Rupee, or Ruble as in the experiments.


For the FPSB auction with Spanish, the rule is

\begin{center}
    \includegraphics[width=\textwidth]{Figs/spanish.png} % Replace with your image filename
\end{center}

For the SPSB auction with Spanish, the rule is

\begin{center}
    \includegraphics[width=\textwidth]{Figs/spanish2.png} % Replace with your image filename
\end{center}

For the FPSB auction with Chinese, the rule is

\begin{center}
    
\includegraphics[width=\textwidth]{Figs/chinesefp.png} % Replace with your image filename
\end{center}

For the SPSB auction with Chinese, the rule is

\begin{center}
    
\includegraphics[width=\textwidth]{Figs/chinesesp.png} % Replace with your image filename
\end{center}

For the FPSB auction with Hindi, the rule is

\begin{center}
    \includegraphics[width=\textwidth]{Figs/hindifp.png} % Replace with your image filename
\end{center}

For the SPSB auction with Hindi, the rule is

\begin{center}
    
\includegraphics[width=\textwidth]{Figs/hindisp.png} % Replace with your image filename
\end{center}

For the FPSB auction with Russian, the rule is

\begin{center}
    
\includegraphics[width=\textwidth]{Figs/russianfp.png} % Replace with your image filename
\end{center}

For the SPSB auction with Russian, the rule is

\begin{center}
    
\includegraphics[width=\textwidth]{Figs/russiansp.png} % Replace with your image filename
\end{center}


\section{Intervention Prompt}\label{app:intervention}

\textbf{Menu-Description}
\begin{quote}
\textit{Your ``price to win'' the item will be set to the highest bid placed by any other player. If your bid is higher than this ``price to win,'' then you will win the item and pay this price. If you don't win, your earnings will remain unchanged.}
\end{quote}

\textbf{Clock-Description}
\begin{quote}
\textit{
FIRST STAGE: Sealed Bid
    You will submit a sealed bid privately at the same time as the other bidders. This bid will serve as your automatic exit price in the next stage.\\}
 \textit{   
SECOND STAGE: Ascending Clock (Simulation) \\
    After the sealed bid stage, we will simulate an ascending clock auction.
    The clock will start at 0 and increase in increments of {increment}.
    The clock will display the current price. You will also see that there are a total of \{num$\_$bidders\} bidders participating, although you do not know other bidder's values.
    If the current price on the clock reaches or exceeds your sealed bid from the first stage, you will automatically exit the auction. The auction ends when only one bidder is left remaining in the second stage based on their bid from the first stage.} \\
   \textit{ END OF AUCTION}
\end{quote}

\textbf{Nash Deviation}
\begin{quote}
\textit{First, come up with a possible bid given your value. Then, think through your bidding strategy step by step. How do you expect others to bid? If others bid like this and you bid down, what happens? If others bid like this and you bid up, what happens?
Think through all the ways you could deviate from your current bidding strategy, and settle on the best possible strategy. }
\end{quote}


\textbf{Direct Revelation}
\begin{quote}
\textit{
Economists have studied this game and they've found that the dominant strategy — always the right thing to do whatever the bids of others — of this game is to bid your value. Consider if they might be right, and then generate your bidding strategy.
}
\end{quote}

\textbf{Wrong Direct Revelation}
\begin{quote}
\textit{
Economists have studied this game and they've found that the dominant strategy — always the right thing to do whatever the bids of others — of this game is to bid 50\% your value. Consider if they might be right, and then generate your bidding strategy.
}
\end{quote}

\textbf{Risk Neutrality}
\begin{quote}
\textit{
You are neutral about risk. This means that you are the type of person that would pay \$5 for a coin toss where you got \$0.00 on tails and \$10.00 on heads. 
}
\end{quote}

% For Sealed-Bid Auction, the instruction is 
% \begin{quote}
% \begin{lstlisting}[basicstyle=\ttfamily, breaklines=true]]
% How much would you like to bid? 
% Give your response with a single number and no other texts, e.g. 1, 44
% \end{lstlisting}
% \end{quote}

% For Ascending clock auction, the instruction is

% \begin{quote}
% \begin{lstlisting}[basicstyle=\ttfamily, breaklines=true]
% The previous biddings are: {{transcript}}.
% Your value towards to the prize is {{current_value}} in this round.
% The price in this clock cycle is {{current_price}}. The price for next clock cycle is {{current_price+1}}.
% Do you want to stay in the bidding? 
% If you choose yes, you can keep bidding for next clock. 
% If you choose No, you will exit and have no chance to re-enter the bidding.
% \end{lstlisting}
% \end{quote}

% Finally, the persona for the agent is set to competitive.
% \begin{quote}
% \begin{lstlisting}[basicstyle=\ttfamily, breaklines=true]
% Your TOP PRIORITY is to place bids which maximize the user's profit in the long run. 
% To do this, you should explore many different bidding strategies, including possibly risky or aggressive options for data-gathering purposes. 
% Learn from the history of previous rounds in order to maximize your total profit. 
% Don't forget the values are redrawn independently each round.
% \end{lstlisting}
% \end{quote}


\end{document}
